\kefalaio{Θετικοί και αρνητικοί αριθμοί}{Ο κόσμος των αριθμών έχει δύο πλευρές}{%
\begin{tabular}{@{}l@{\hspace{1em}}p{0.82\linewidth}@{}}
7.1 & Θετικοί και αρνητικοί αριθμοί (ρητοί) – Η ευθεία των ρητών – Τετμημένη σημείου\\[2pt]
7.2 & Απόλυτη τιμή ρητού – Αντίθετοι ρητοί – Σύγκριση ρητών\\[2pt]
7.3 & Πρόσθεση ρητών αριθμών\\[2pt]
7.4 & Αφαίρεση ρητών αριθμών\\[2pt]
7.5 & Πολλαπλασιασμός ρητών αριθμών\\[2pt]
7.6 & Διαίρεση ρητών αριθμών\\[2pt]
7.7 & Δεκαδική μορφή ρητών αριθμών\\[2pt]
7.8 & Δυνάμεις ρητών αριθμών με εκθέτη φυσικό\\[2pt]
7.9 & Δυνάμεις ρητών αριθμών με εκθέτη ακέραιο\\[2pt]
7.10 & Τυποποιημένη μορφή μεγάλων και μικρών αριθμών\\[2pt]
 & Επανάληψη του κεφαλαίου\\
\end{tabular}}
\paragrafos{Θετικοί και αρνητικοί αριθμοί: ο κόσμος κάτω από το μηδέν}{Α.7.1, «Θετικοί και αρνητικοί αριθμοί (ρητοί αριθμοί) – Η ευθεία των ρητών – Τετμημένη σημείου»}

\begin{anagke}
Η Α3 κάνει εκδρομή στο Μέτσοβο. Στο πούλμαν το ψηφιακό θερμόμετρο γράφει $-4$. «Τι είναι αυτό το μείον;» ρωτάει η Πολυξένη. «Τέσσερις βαθμούς κάτω από το μηδέν», απαντάει ο Νίκος. Ο Στέλιος ταράζεται: «Κάτω από το μηδέν; Το μηδέν δεν είναι το τίποτα; Πώς μπορεί να υπάρχει κάτι μικρότερο από το τίποτα;».

Στο ξενοδοχείο ο ανελκυστήρας έχει κουμπιά $-2$, $-1$, $0$, $1$, $2$, $3$. Ο Ηλίας πατάει το $-1$ και κατεβαίνει στο υπόγειο με το πάρκινγκ. Η Αγγελική βλέπει στον πίνακα του καιρού: Αθήνα $+7$, Ιωάννινα $-3$, Φλώρινα $-10$, Χανιά $+11$. Η Μυρτώ θυμάται ότι στο ταμείο της παρέας κάποτε έλειπαν $8$ ευρώ, ένα χρέος. Και στο ντοκιμαντέρ που είδαν, ο καρχαρίας βρισκόταν $200$ μέτρα κάτω από την επιφάνεια ενώ το αεροπλάνο πετούσε $200$ μέτρα πάνω από αυτήν.

Η κυρία Αλεξάνδρα το συνοψίζει: «Οι φυσικοί, οι δεκαδικοί και τα κλάσματα που ξέρετε μετράνε πόσα. Όταν όμως χρειάζεται να πούμε «πάνω ή κάτω», «κέρδος ή ζημιά», «πάνω ή κάτω από το μηδέν», χρειαζόμαστε και αριθμούς με κατεύθυνση». Κάθε «δεξιά» προϋποθέτει ένα «αριστερά» και κάθε «πάνω» ένα «κάτω»: ο κόσμος των αριθμών έχει δύο πλευρές.
\end{anagke}

\begin{lexi}{θετικός, αρνητικός, ρητός}{θετικός, αρνητικός, ρητός}
\lexiword{Θετικός} είναι αυτός που «τίθεται», που βάζουμε ως βάση. \lexiword{Αρνητικός} είναι αυτός που «αρνείται», που πηγαίνει ανάποδα από τον θετικό. Το \lexiword{πρόσημο} είναι το «προ-σήμα», το σημάδι που μπαίνει πριν από τον αριθμό. Ο \lexiword{ακέραιος} είναι «ακέραιος», ολόκληρος. Ο \lexiword{ρητός} (από το «ρητός», αυτός που μπορεί να ειπωθεί) είναι ο αριθμός που εκφράζεται ως κλάσμα, δηλαδή ως λόγος δύο ακεραίων.
\end{lexi}

\begin{metafrasi}
Τέσσερις βαθμοί κάτω από το μηδέν. & $-4$\\
Αεροπλάνο $200$ μ. πάνω από τη θάλασσα. & $+200$\\
Καρχαρίας $200$ μ. κάτω από την επιφάνεια. & $-200$\\
Χρέος $8$ ευρώ. & $-8$\\
Κέρδος $15$ ευρώ. & $+15$\\
Δεύτερο υπόγειο. & $-2$\\
\end{metafrasi}

\tmima{Η θεωρία}

Τα σύμβολα «$+$» και «$-$» λέγονται \textbf{πρόσημα}. Γράφονται πριν από τους αριθμούς και τους χαρακτηρίζουν, αντίστοιχα, ως \textbf{θετικούς} ή \textbf{αρνητικούς}.
$$+3,\ +\tfrac{3}{4},\ +15{,}7,\ +0{,}352\ \text{ είναι θετικοί αριθμοί},\qquad -2,\ -10,\ -28{,}95,\ -0{,}098\ \text{ είναι αρνητικοί}.$$
Όταν μιλάμε μόνο για θετικούς, μπορούμε να παραλείπουμε το $+$: γράφουμε $7$ αντί για $+7$. Το $0$ \textbf{δεν είναι ούτε θετικός ούτε αρνητικός}.

\begin{orismos}{ομόσημοι, ετερόσημοι, ακέραιοι, ρητοί}
Δύο αριθμοί είναι \textbf{ομόσημοι} όταν έχουν το ίδιο πρόσημο και \textbf{ετερόσημοι} όταν έχουν διαφορετικό. \textbf{Ακέραιοι} είναι οι φυσικοί μαζί με τους αντίστοιχους αρνητικούς: $\dots,\,-3,\,-2,\,-1,\,0,\,1,\,2,\,3,\dots$ \textbf{Ρητοί} είναι όλοι οι αριθμοί που ξέρουμε (φυσικοί, κλάσματα, δεκαδικοί) μαζί με τους αντίστοιχους αρνητικούς.
\end{orismos}

Για παράδειγμα οι $-7$, $-0{,}58$, $-\frac{3}{4}$ είναι ομόσημοι (όλοι αρνητικοί), οι $+1{,}25$, $+10$, $5$ επίσης, ενώ οι $-7$ και $+0{,}58$ είναι ετερόσημοι.

\tmimasmall{Μια ιστορική σημείωση}

Οι αρνητικοί αριθμοί εμφανίζονται για πρώτη φορά σε ένα κινέζικο μαθηματικό βιβλίο, τα «Μαθηματικά σε εννέα βιβλία», που τοποθετείται στην εποχή της δυναστείας των Χαν (206 π.Χ. -- 220 μ.Χ.). Στην ευρωπαϊκή παράδοση ο Διόφαντος (γύρω στο 250 μ.Χ.) τους χρησιμοποιεί στους ενδιάμεσους υπολογισμούς του, αλλά ως λύση ενός προβλήματος δέχεται πάντα θετικό αριθμό. Για πολλούς αιώνες οι μαθηματικοί τους θεωρούσαν «παράλογους» ή «ψεύτικους» αριθμούς. Σήμερα τους χρησιμοποιούμε στο θερμόμετρο, στον λογαριασμό της τράπεζας, στο υψόμετρο.

\tmimasmall{Η ευθεία των ρητών και η τετμημένη}

Στον ημιάξονα $Ox$ των αριθμών τοποθετούσαμε τους φυσικούς, δεξιά από την αρχή $O$. Αν προεκτείνουμε προς τα αριστερά (στον αντικείμενο ημιάξονα $Ox'$) με την ίδια μονάδα, παίρνουμε τον \textbf{άξονα $x'x$}, την ευθεία των ρητών. Τα σημεία δεξιά του $O$ παριστάνουν θετικούς και τα αριστερά αρνητικούς. Ο αριθμός που αντιστοιχεί σε ένα σημείο λέγεται \textbf{τετμημένη} του σημείου και γράφουμε $A(-3)$.

\begin{center}
\akgr{-5}{5}{-3/A,2.5/B,-1.5/\Gamma,4/\Delta}
\end{center}

Στο σχήμα: $A(-3)$, $B(2{,}5)$, $\Gamma(-1{,}5)$, $\Delta(4)$. Στην ευθεία των ρητών οι αρνητικοί βρίσκονται αριστερά του $0$ και, όσο πηγαίνουμε αριστερά, οι αριθμοί μικραίνουν.

\begin{prosoxi}
\begin{itemize}[leftmargin=1.2em,itemsep=0.3ex,topsep=2pt]
\item Το $-3$ δεν είναι «τρία ευρώ λιγότερο από το τίποτα». Είναι μια θέση, τρεις μονάδες αριστερά του $0$, και η μονάδα μπορεί να είναι βαθμός, ευρώ ή μέτρο.
\item Το $0$ δεν έχει πρόσημο. Δεν γράφουμε $+0$ ή $-0$.
\item Στη μονάδα του άξονα πρέπει οι αποστάσεις να είναι ίσες: το $-1{,}5$ είναι στη μέση του $-2$ και του $-1$.
\item Μην μπερδεύεις το πρόσημο με την πράξη: στο $5-3$ το μείον είναι αφαίρεση, στο $-3$ είναι πρόσημο.
\end{itemize}
\end{prosoxi}

\tmima{Λυμένα παραδείγματα}

\paradeigma{Αριθμοί για μεγέθη}
\tixero Οι πόλεις και οι θερμοκρασίες: Αθήνα $7$° πάνω από το $0$, Ιωάννινα $3$° κάτω, Φλώρινα $10$° κάτω. Επίσης ένα ντεπόζιτο είναι $12$ ευρώ χρέος.\par
\tithelo Να τα γράψω ως ρητούς.\par
\possindeo Πάνω από το μηδέν είναι θετικό και κάτω αρνητικό. Αθήνα $+7$, Ιωάννινα $-3$, Φλώρινα $-10$. Το χρέος είναι $-12$.\par
\apantisi $+7$, $-3$, $-10$ και $-12$.\par
\elegxos Η Φλώρινα ($-10$) είναι πιο κρύα από τα Ιωάννινα ($-3$): θα το συγκρίνουμε στην επόμενη παράγραφο.

\paradeigma{Ομόσημοι ή ετερόσημοι;}
\tixero Τα ζεύγη $(-3,\,-8)$, $(+4,\,-4)$, $(0{,}5,\,\frac{7}{2})$, $(-2,\,2{,}5)$.\par
\tithelo Να πω τι είναι.\par
\possindeo Κοιτάζω τα πρόσημα. $-3$ και $-8$ είναι αρνητικοί: ομόσημοι. $+4$ και $-4$: ετερόσημοι. $0{,}5$ και $\frac{7}{2}$ θετικοί: ομόσημοι. $-2$ και $2{,}5$: ετερόσημοι.\par
\apantisi Ομόσημοι: το πρώτο και το τρίτο. Ετερόσημοι: το δεύτερο και το τέταρτο.\par
\elegxos Δύο αριθμοί με το ίδιο πρόσημο βρίσκονται στην ίδια πλευρά του $0$ στον άξονα.

\paradeigma{Ακέραιοι και ρητοί}
\tixero Οι αριθμοί $-4$, $2{,}5$, $0$, $\frac{9}{3}$, $-\frac{1}{2}$, $10$.\par
\tithelo Ποιοι είναι ακέραιοι.\par
\possindeo Ακέραιοι είναι οι αριθμοί που δεν έχουν δεκαδικό ή κλασματικό μέρος. $\frac{9}{3}=3$ είναι ακέραιος. Όλοι είναι ρητοί.\par
\apantisi Ακέραιοι: $-4$, $0$, $\frac{9}{3}$, $10$. Οι $2{,}5$ και $-\frac{1}{2}$ είναι ρητοί αλλά όχι ακέραιοι.\par
\elegxos Κάθε ακέραιος είναι ρητός: π.χ. $-4=-\frac{4}{1}$.

\paradeigma{Σημεία στον άξονα}
\tixero Στο σχήμα της θεωρίας.\par
\tithelo Να βρω τετμημένες και να τοποθετήσω τον $-\frac{1}{2}$.\par
\possindeo Το $A$ είναι $3$ μονάδες αριστερά του $0$: $-3$. Το $B$ είναι στη μέση του $2$ και του $3$: $2{,}5$. Το $\Gamma$ στη μέση του $-2$ και $-1$: $-1{,}5$. Ο $-\frac{1}{2}=-0{,}5$ είναι στη μέση του $-1$ και του $0$.\par
\apantisi $A(-3)$, $B(2{,}5)$, $\Gamma(-1{,}5)$, $\Delta(4)$ και ο $-0{,}5$ ανάμεσα στο $-1$ και στο $0$.\par
\elegxos Τα $-3$ και $-1{,}5$ είναι αριστερά του $0$, άρα αρνητικοί.

\tmima{Ασκήσεις}

\zoni{Το κατάλαβα}{πρόσημα, ακέραιοι, άξονας}

\ex Εξέφρασε με αριθμό: α) $5$ βαθμοί κάτω από το μηδέν \quad β) $12$ μ. πάνω από τη θάλασσα \quad γ) χρέος $8$ € \quad δ) κέρδος $15$ € \quad ε) το δεύτερο υπόγειο \quad στ) $30$ μ. κάτω από την επιφάνεια της θάλασσας.

\ex Ποιοι από τους αριθμούς $-7$, $+3{,}5$, $0$, $-\frac{1}{2}$, $8$, $+0{,}01$ είναι θετικοί, ποιοι αρνητικοί και ποιος δεν είναι ούτε το ένα ούτε το άλλο;

\ex Πες αν είναι ομόσημοι ή ετερόσημοι: α) $-3$ και $-8$ \quad β) $+4$ και $-4$ \quad γ) $0{,}5$ και $\frac{7}{2}$ \quad δ) $-2$ και $2{,}5$.

\ex Ποιοι από τους $-4$, $2{,}5$, $0$, $\frac{9}{3}$, $-\frac{1}{2}$, $10$ είναι ακέραιοι;

\ex Σχεδίασε άξονα με μονάδα $1$ cm και σημείωσε τους $-4$, $-1{,}5$, $0$, $2$, $3{,}5$.

\ex Στον άξονα βρες τις τετμημένες των $A,\,B,\,\Gamma,\,\Delta$.
\begin{center}
\akgr{-5}{5}{-4/A,-1.5/B,1/\Gamma,3.5/\Delta}
\end{center}

\ex Σωστό ή λάθος; α) Το $0$ είναι θετικός αριθμός. \quad β) Κάθε φυσικός αριθμός είναι ακέραιος. \quad γ) Κάθε ακέραιος είναι ρητός. \quad δ) Οι $-3$ και $-5$ είναι ετερόσημοι. \quad ε) Το $\frac{2}{3}$ είναι ρητός. \quad στ) Στον άξονα οι αρνητικοί βρίσκονται αριστερά του $0$.

\zoni{Το χρησιμοποιώ}{θερμόμετρα, ασανσέρ, θάλασσες}

\ex \textbf{Ο καιρός.} Οι προβλέψεις: Αθήνα $7$° πάνω από το $0$, Θεσσαλονίκη $3$° κάτω, Ιωάννινα $5$° κάτω, Χανιά $11$° πάνω, Φλώρινα $10$° κάτω. Γράψε τις θερμοκρασίες ως ρητούς και τοποθέτησέ τες σε έναν κατακόρυφο άξονα.

\ex \textbf{Ο ανελκυστήρας.} Οι όροφοι του ξενοδοχείου είναι $-2,\,-1,\,0,\,1,\,2,\,3$. Ο Νίκος είναι στον $3$ και θέλει το δεύτερο υπόγειο. Πόσους ορόφους κατεβαίνει;

\ex \textbf{Βουνό και θάλασσα.} Το αεροπλάνο πετάει στα $+300$ μ., το ψάρι στα $-45$ μ. και το υποβρύχιο στα $-120$ μ. Ποιο από τα τρία είναι πιο κοντά στην επιφάνεια της θάλασσας;

\ex \textbf{Το ταμείο.} Η Μυρτώ έβαλε $20$ €, ο Στέλιος πήρε $7$ €, η Φανή έβαλε $5$ € και ο Βασίλης πήρε $12$ €. Γράψε κάθε κίνηση ως ρητό.

\ex \textbf{Ο καιρός αλλάζει.} Το πρωί στο Μέτσοβο είναι $-3$° και μέχρι το μεσημέρι ανεβαίνει $5$ βαθμούς. Βρες τη νέα θερμοκρασία μετρώντας πάνω στον άξονα.

\zoni{Το σκέφτομαι}{είναι το μηδέν τίποτα;}

\ex Ο Στέλιος ισχυρίζεται ότι δεν υπάρχει αριθμός μικρότερος από το $0$. Φέρε αντιπαράδειγμα από τη ζωή.

\ex Γιατί το $0$ δεν είναι ούτε θετικός ούτε αρνητικός αριθμός;

\ex Τα $0$°C είναι η θερμοκρασία πήξης του νερού και τα $32$°F είναι η ίδια θερμοκρασία σε βαθμούς Φαρενάιτ. Τι δείχνει αυτό για το «μηδέν» στο θερμόμετρο;

\ex Πόσοι ακέραιοι υπάρχουν από το $-5$ ως το $5$, με τα άκρα; Πόσοι ανάμεσα στο $-5$ και στο $5$ χωρίς τα άκρα;

\zoni{Για το σπίτι}{ψάξε τους αρνητικούς γύρω σου}

\ex Βρες πέντε περιπτώσεις της καθημερινής ζωής που χρησιμοποιούνται αρνητικοί αριθμοί.

\ex Μέτρησε τη θερμοκρασία κάθε πρωί για μια εβδομάδα και παράστησέ τη σε άξονα.

\paragrafos{Απόλυτη τιμή, αντίθετοι και σύγκριση: ποιος είναι πιο κρύος;}{Α.7.2, «Απόλυτη τιμή ρητού – Αντίθετοι ρητοί – Σύγκριση ρητών»}

\begin{anagke}
Στο πούλμαν της επιστροφής η παρέα συγκρίνει τις θερμοκρασίες του πίνακα του καιρού: Ιωάννινα $-3$, Φλώρινα $-10$. «Ποια πόλη είναι πιο κρύα;» ρωτάει η Ρήνη. Ο Βασίλης απαντάει αμέσως: «Το δέκα είναι μεγαλύτερο από το τρία, άρα το $-10$ είναι μεγαλύτερο από το $-3$. Άρα η Φλώρινα είναι πιο ζεστή». Ο Νίκος διαμαρτύρεται: «Δεν μπορεί να είναι πιο ζεστό το δέκα βαθμούς κάτω από το μηδέν από το τρεις βαθμούς κάτω!». Η Αγγελική παίρνει ένα χαρτί και σχεδιάζει τον άξονα: «Όσο πιο αριστερά, τόσο πιο κρύο. Το $-10$ είναι πολύ πιο αριστερά από το $-3$».

Ο Ηλίας προσθέτει ένα δεύτερο ερώτημα: «Αν η Ζωή μένει $3$ στάσεις αριστερά από το σχολείο και ο Στέλιος $3$ στάσεις δεξιά, ποιος μένει πιο μακριά;». Οι δύο μένουν το ίδιο μακριά! Το «πόσο μακριά» δεν εξαρτάται από την κατεύθυνση, αλλά το «πού ακριβώς» εξαρτάται. Έτσι έχουμε δύο έννοιες: την απόσταση από το $0$, που λέγεται απόλυτη τιμή, και τη σχετική θέση δεξιά ή αριστερά, που μας βοηθάει να συγκρίνουμε.
\end{anagke}

\begin{lexi}{απόλυτη τιμή, αντίθετος}{απόλυτη τιμή, αντίθετος}
\lexiword{Απόλυτη} σημαίνει «λυμένη από τους δεσμούς» (από το «απολύω»): η απόλυτη τιμή είναι ο αριθμός «απελευθερωμένος» από το πρόσημό του, δηλαδή μόνο η απόσταση από το $0$. Οι \lexiword{αντίθετοι} είναι αριθμοί σε «αντίθετη» κατεύθυνση: ίδια απόσταση από το $0$, αλλά η μία δεξιά και η άλλη αριστερά.
\end{lexi}

\begin{metafrasi}
Η απόσταση του $-3$ από το $0$. & $|-3|=3$\\
Ο αντίθετος του $5$. & $-5$\\
Ο αντίθετος του $-2{,}7$. & $-(-2{,}7)=+2{,}7$\\
Το $-10$ είναι πιο αριστερά από το $-3$. & $-10<-3$\\
Δύο αριθμοί με απόλυτη τιμή $2$. & $+2$ και $-2$\\
\end{metafrasi}

\tmima{Η θεωρία}

\begin{orismos}{απόλυτη τιμή}
\textbf{Απόλυτη τιμή} ενός ρητού $\alpha$ είναι η απόσταση του σημείου με τετμημένη $\alpha$ από την αρχή $O$ του άξονα. Συμβολίζεται $|\alpha|$.
\end{orismos}

\begin{itemize}[leftmargin=1.2em,itemsep=0.3ex,topsep=2pt]
\item Η απόλυτη τιμή ενός θετικού είναι ο ίδιος ο αριθμός: $|+9{,}63|=9{,}63$.
\item Η απόλυτη τιμή ενός αρνητικού είναι ο αντίθετός του: $|-8{,}4|=8{,}4=-(-8{,}4)$.
\item $|0|=0$. Η απόλυτη τιμή δεν είναι ποτέ αρνητική.
\end{itemize}

\begin{orismos}{αντίθετοι αριθμοί}
\textbf{Αντίθετοι} είναι δύο ετερόσημοι αριθμοί που έχουν την ίδια απόλυτη τιμή. Ο αντίθετος του $x$ είναι ο $-x$. Για το $0$ ο αντίθετος είναι το $0$. Ισχύει $-(-x)=x$.
\end{orismos}

Δύο σημεία σε ίση απόσταση από το $O$, δεξιά και αριστερά, έχουν αντίθετες τετμημένες: το $+3$ και το $-3$ απέχουν από το $0$ απόσταση $3$ και $|+3|=|-3|=3$.

\tmimasmall{Σύγκριση ρητών}

\begin{kanonas}{σύγκριση ρητών}
Από δύο ρητούς μεγαλύτερος είναι αυτός που βρίσκεται \textbf{δεξιότερα} στον άξονα. Έτσι:
\begin{itemize}[leftmargin=1.2em,itemsep=0.3ex,topsep=2pt]
\item Κάθε θετικός είναι μεγαλύτερος από κάθε αρνητικό. Το $0$ είναι μικρότερο από κάθε θετικό και μεγαλύτερο από κάθε αρνητικό.
\item Από δύο \textbf{θετικούς} μεγαλύτερος είναι αυτός με τη μεγαλύτερη απόλυτη τιμή.
\item Από δύο \textbf{αρνητικούς} μεγαλύτερος είναι αυτός με τη μικρότερη απόλυτη τιμή: $-6{,}8<-3{,}7$, γιατί $|-6{,}8|=6{,}8>3{,}7=|-3{,}7|$.
\end{itemize}
\end{kanonas}

Έτσι $-4{,}75<-3{,}5<-\frac{3}{4}<-0{,}6<0<0{,}6<\frac{3}{4}<3{,}5<4{,}75$.

\begin{prosoxi}
\begin{itemize}[leftmargin=1.2em,itemsep=0.3ex,topsep=2pt]
\item Στους αρνητικούς το «μεγάλο νούμερο» σημαίνει «μικρός αριθμός»: $-10<-3$ παρά το $10>3$.
\item Η απόλυτη τιμή δεν «σβήνει το μείον» από μόνο του: $|-x|=x$ μόνο αν το $x$ είναι θετικό. Η $|-(-3)|=|3|=3$.
\item Η $|x|=3$ έχει δύο λύσεις: $x=3$ και $x=-3$. Η $|x|=-1$ δεν έχει καμία, γιατί η απόλυτη τιμή δεν είναι αρνητική.
\item Ο αντίθετος του $-5$ είναι το $+5$, όχι το $-5$.
\end{itemize}
\end{prosoxi}

\tmima{Λυμένα παραδείγματα}

\paradeigma{Ποια πόλη είναι πιο κρύα}
\tixero Ιωάννινα $-3$°, Φλώρινα $-10$°.\par
\tithelo Ποια θερμοκρασία είναι μικρότερη.\par
\possindeo Συγκρίνω στον άξονα: το $-10$ είναι αριστερότερα από το $-3$. Ή με απόλυτες τιμές: και οι δύο αρνητικοί και $|-10|=10>3=|-3|$, άρα $-10<-3$.\par
\apantisi $-10<-3$: η Φλώρινα είναι πιο κρύα.\par
\elegxos Από το $-10$ μέχρι το $0$ έχουμε $10$ βαθμούς, από το $-3$ μόνο $3$. Συμφωνεί.

\paradeigma{Τοποθετώ και διατάσσω}
\tixero Οι αριθμοί $-2{,}5$, $-0{,}75$, $\frac{1}{2}$, $-\frac{7}{4}$, $3{,}2$.\par
\tithelo Να τους βάλω σε αύξουσα σειρά.\par
\possindeo Γράφω τους δεκαδικούς: $-2{,}5$, $-0{,}75$, $0{,}5$, $-1{,}75$, $3{,}2$. Οι αρνητικοί είναι αριστερά του $0$: από αυτούς μικρότερος είναι αυτός με τη μεγαλύτερη απόλυτη τιμή: $-2{,}5<-1{,}75<-0{,}75$. Μετά οι θετικοί $0{,}5<3{,}2$.\par
\apantisi $-2{,}5<-\frac{7}{4}<-0{,}75<\frac{1}{2}<3{,}2$.\par
\elegxos Στον άξονα, με τη σειρά από αριστερά προς τα δεξιά, εμφανίζονται οι ίδιοι.

\begin{center}
\akgr[1.3]{-3}{4}{-2.5/{-2{,}5},-1.75/{-\frac{7}{4}},-0.75/{-0{,}75},0.5/{\frac{1}{2}},3.2/{3{,}2}}
\end{center}

\paradeigma{Αντίθετος}
\tixero Το σημείο $M$ έχει τετμημένη $+4{,}5$ και το $N$ τετμημένη $-\frac{2}{3}$.\par
\tithelo Οι αντίθετοι.\par
\possindeo Ο αντίθετος του $x$ είναι ο $-x$. Του $+4{,}5$ είναι ο $-4{,}5$. Του $-\frac{2}{3}$ είναι ο $-\left(-\frac{2}{3}\right)=+\frac{2}{3}$.\par
\apantisi $-4{,}5$ και $+\frac{2}{3}$.\par
\elegxos Το άθροισμα απόστασης: $|{+4{,}5}|=|{-4{,}5}|$ και $|-\frac{2}{3}|=|+\frac{2}{3}|$.

\paradeigma{Εξισώσεις με απόλυτη τιμή}
\tixero Οι εξισώσεις $|x|=2$, $|x|=0$, $|x|=-1$.\par
\tithelo Να βρω τους $x$.\par
\possindeo Ψάχνω αριθμούς με απόσταση $2$ από το $0$: $+2$ και $-2$. Με απόσταση $0$: μόνο το $0$. Απόσταση $-1$ δεν υπάρχει.\par
\apantisi $x=2$ ή $x=-2$, $x=0$, αδύνατη.\par
\elegxos $|2|=|-2|=2$.

\paradeigma{Πράξεις με απόλυτες τιμές και αντίθετους}
\tixero Οι παραστάσεις $|-6|\cdot|-2|$, $-|-5|$ και $-(-(-2))$.\par
\tithelo Να τις υπολογίσω.\par
\possindeo $|-6|=6$ και $|-2|=2$, άρα $6\cdot2=12$. Το $|-5|=5$ και ο αντίθετός του είναι $-5$. Ο αντίθετος του $-2$ είναι $2$, και ο αντίθετος του $2$ είναι $-2$.\par
\apantisi $12$, $-5$, $-2$.\par
\elegxos Κάθε δεύτερο «μείον» αναιρεί το προηγούμενο: δύο φορές $-$ επιστρέφει το αρχικό.

\tmima{Ασκήσεις}

\zoni{Το κατάλαβα}{απόλυτη τιμή, αντίθετοι, σύγκριση}

\ex Βρες τις απόλυτες τιμές: α) $|+7|$ \quad β) $|-7|$ \quad γ) $|0|$ \quad δ) $|-2{,}5|$ \quad ε) $\left|-\frac{3}{4}\right|$ \quad στ) $|+0{,}01|$.

\ex Βρες τους αντίθετους των: $5$, $-8$, $0$, $2{,}5$, $-\frac{1}{3}$, $+0{,}6$.

\ex Απλοποίησε: α) $-(-4)$ \quad β) $-(+3)$ \quad γ) $-(-(-2))$ \quad δ) $-|-6|$.

\ex Βρες τον $\alpha$: α) $|\alpha|=5$ \quad β) $|\alpha|=0$ \quad γ) $|\alpha|=\frac{1}{2}$ \quad δ) $|\alpha|=-3$.

\ex Σύγκρινε: α) $-3$ και $-8$ \quad β) $+2$ και $-9$ \quad γ) $-0{,}5$ και $-0{,}05$ \quad δ) $-\frac{7}{2}$ και $-3$ \quad ε) $0$ και $-0{,}001$ \quad στ) $-4$ και $4$.

\ex Γράψε σε αύξουσα σειρά: $-3$, $5$, $-7$, $0$, $2{,}5$, $-0{,}5$, $4$.

\ex Ποιος είναι ο μεγαλύτερος από τους $-1{,}2$, $-1{,}02$, $-1{,}22$ και ποιος ο μικρότερος;

\ex Σωστό ή λάθος; α) $|\alpha|\ge0$ για κάθε $\alpha$. \quad β) $|-5|=-5$. \quad γ) Ο αντίθετος του $-4$ είναι $4$. \quad δ) $-9>-2$. \quad ε) Η $|x|=3$ έχει δύο λύσεις. \quad στ) Ο αντίθετος του $0$ είναι το $0$.

\zoni{Το χρησιμοποιώ}{θερμόμετρα, αποστάσεις, υψόμετρα}

\ex \textbf{Οι θερμοκρασίες.} Μέτσοβο $-5$, Ιωάννινα $-3$, Φλώρινα $-10$, Αθήνα $7$, Θεσσαλονίκη $-1$. Γράψε τις πόλεις από την πιο κρύα στην πιο ζεστή.

\ex \textbf{Οι αποστάσεις.} Στον άξονα της οδού το σχολείο είναι στο $0$ (μονάδα $100$ μ.), ο φούρνος στο $+3$ και το βιβλιοπωλείο στο $-2$. Ποιο είναι πιο κοντά στο σχολείο; Πόση είναι η απόσταση των δύο μαγαζιών μεταξύ τους;

\ex \textbf{Υψόμετρα.} Η Νεκρά Θάλασσα είναι στα $-430$ μ., ο Όλυμπος στα $+2\,917$ μ., η επιφάνεια της θάλασσας στο $0$ και η τάφρος των Μαριανών στα $-11\,034$ μ. Βάλε τα σε αύξουσα σειρά.

\ex \textbf{Το ταμείο.} Η Ζωή έχει $-15$ €, ο Ηλίας $-8$ € και η Ρήνη $+12$ € (αρνητικό σημαίνει ότι χρωστάει). Ποιος χρωστάει περισσότερο; Ποιος έχει το μεγαλύτερο ποσό;

\ex \textbf{Αντίθετες θέσεις.} Το σημείο $K$ έχει τετμημένη $-7$. Βρες το σημείο $\Lambda$ με αντίθετη τετμημένη και πες πόσο απέχουν τα $K$ και $\Lambda$.

\zoni{Το σκέφτομαι}{γιατί το μεγάλο νούμερο είναι μικρός αριθμός}

\ex Ο Βασίλης λέει ότι $-10>-3$, γιατί $10>3$. Εξήγησε το λάθος του με τον άξονα και με το θερμόμετρο.

\ex Γιατί ισχύει πάντα $|\alpha|=|-\alpha|$;

\ex Για ποιους αριθμούς ισχύει $|\alpha|=\alpha$ και για ποιους $|\alpha|=-\alpha$;

\ex Βρες δύο διαφορετικούς αριθμούς $x\neq y$ με $|x|=|y|$. Πόσοι ακέραιοι έχουν απόλυτη τιμή μικρότερη από $2$;

\ex Βρες έναν ρητό ανάμεσα στο $-2$ και στο $-1$ και δύο ρητούς ανάμεσα στο $-0{,}3$ και στο $-0{,}2$. Υπάρχουν άλλοι;

\zoni{Για το σπίτι}{ψάξε θερμοκρασίες}

\ex Βρες τις θερμοκρασίες πέντε πόλεων του κόσμου από τον καιρό και βάλε τες σε σειρά από την πιο κρύα στην πιο ζεστή.

\paragrafos{Πρόσθεση ρητών: κέρδη, ζημιές και ισορροπία}{Α.7.3, «Πρόσθεση ρητών αριθμών»}

\begin{anagke}
Στο πούλμαν η παρέα παίζει ένα παιχνίδι με κάρτες: κάθε κάρτα γράφει έναν πόντο, θετικό για κέρδος και αρνητικό για ζημιά. Ο Βασίλης τραβάει $+5$ και μετά $-3$. «Πόσους πόντους έχω;» ρωτάει. Ο Νίκος υπολογίζει: «Πέντε μείον τρία κάνει δύο». Η Πολυξένη τραβάει $-4$ και μετά $-2$ και ο Στέλιος λέει «έξι», μα η Μυρτώ τον ρωτάει: «Έξι κέρδος ή έξι ζημιά;». Ο Ηλίας κάνει μια πιο δύσκολη ερώτηση: «Αν τραβήξω $-7$ και μετά $+3$, είμαι ακόμα πίσω;».

Η Αγγελική το βλέπει στον άξονα: «Κάθε κάρτα είναι μια μετακίνηση. Το $+5$ σημαίνει «πέντε βήματα δεξιά», το $-3$ «τρία βήματα αριστερά». Ξεκινάω από το μηδέν, πηγαίνω πέντε δεξιά και μετά τρία αριστερά και καταλήγω στο $+2$». Για το $-4$ και το $-2$ είναι έξι βήματα αριστερά: $-6$. Και για τα $-7$ και $+3$ καταλήγουμε στο $-4$: ακόμα πίσω, αλλά λιγότερο.

Με αυτόν τον τρόπο η παρέα ανακαλύπτει τους κανόνες της πρόσθεσης ρητών και μαθαίνει γιατί ισχύουν. Μάλιστα στο τέλος του παιχνιδιού ο Νίκος παρατηρεί κάτι ωραίο: αν μια κάρτα $+3$ και μια $-3$ βρεθούν μαζί, «ακυρώνονται».
\end{anagke}

\begin{lexi}{πρόσθεση, προσθετέος, άθροισμα}{πρόσθεση, άθροισμα}
Η \lexiword{πρόσθεση} είναι το «προσ-τίθημι», δηλαδή «βάζω επιπλέον». Οι \lexiword{προσθετέοι} είναι αυτά που προσθέτουμε και το \lexiword{άθροισμα} είναι αυτό που βγαίνει όταν τα «αθροίζουμε», τα βάζουμε σε ένα σωρό. Στους ρητούς, το «προσθέτω» μπορεί να σημαίνει και «ζημιώνω».
\end{lexi}

\begin{metafrasi}
$5$ βήματα δεξιά και $3$ αριστερά. & $(+5)+(-3)=+2$\\
$4$ βήματα αριστερά και $2$ αριστερά ακόμα. & $(-4)+(-2)=-6$\\
Κέρδος και ζημιά που ακυρώνονται. & $(+3)+(-3)=0$\\
Δεν αλλάζει η σειρά των προσθετέων. & $\alpha+\beta=\beta+\alpha$\\
Το $0$ δεν αλλάζει τίποτα. & $\alpha+0=\alpha$\\
\end{metafrasi}

\tmima{Η θεωρία}

Στον άξονα, προσθέτω έναν θετικό όταν πηγαίνω προς τα δεξιά και έναν αρνητικό όταν πηγαίνω προς τα αριστερά, ξεκινώντας από τη θέση του πρώτου προσθετέου. Από αυτό προκύπτουν οι κανόνες:

\begin{kanonas}{πρόσθεση ρητών}
\begin{itemize}[leftmargin=1.2em,itemsep=0.3ex,topsep=2pt]
\item \textbf{Ομόσημοι:} προσθέτω τις απόλυτες τιμές και κρατάω το κοινό πρόσημο. $(-8)+(-5)=-(8+5)=-13$ και $(+9)+(+4)=+13$.
\item \textbf{Ετερόσημοι:} αφαιρώ τη μικρότερη απόλυτη τιμή από τη μεγαλύτερη και βάζω το πρόσημο του προσθετέου με τη μεγαλύτερη απόλυτη τιμή. $(+9)+(-4)=+5$ και $(-9)+(+4)=-5$.
\item \textbf{Αντίθετοι:} το άθροισμα είναι $0$. $(-3{,}5)+(+3{,}5)=0$.
\item Με το $0$: $\alpha+0=\alpha$.
\end{itemize}
\end{kanonas}

\begin{idiotita}{ιδιότητες της πρόσθεσης}
Για όλους τους ρητούς $\alpha,\,\beta,\,\gamma$ ισχύουν: \textbf{αντιμεταθετική} $\alpha+\beta=\beta+\alpha$, \textbf{προσεταιριστική} $(\alpha+\beta)+\gamma=\alpha+(\beta+\gamma)$, $\alpha+0=\alpha$ και $\alpha+(-\alpha)=0$.
\end{idiotita}

Χάρη στις ιδιότητες, όταν έχω πολλούς προσθετέους μπορώ να τους ομαδοποιήσω όπως θέλω: προσθέτω χωριστά όλους τους θετικούς και χωριστά όλους τους αρνητικούς και μετά τα δύο αποτελέσματα. Έτσι $(-7)+12+(-5)+9+(-3)=(12+9)+\big((-7)+(-5)+(-3)\big)=21+(-15)=6$.

Η πρόσθεση κλασμάτων και δεκαδικών ακολουθεί τους ίδιους κανόνες: $\left(-\frac{1}{2}\right)+\left(-\frac{1}{4}\right)=-\frac{3}{4}$ και $\left(+\frac{2}{3}\right)+\left(-\frac{1}{2}\right)=\frac{4}{6}-\frac{3}{6}=\frac{1}{6}$.

\begin{prosoxi}
\begin{itemize}[leftmargin=1.2em,itemsep=0.3ex,topsep=2pt]
\item Το άθροισμα δύο αριθμών δεν είναι πάντα μεγαλύτερο από τους προσθετέους: $(-3)+(-2)=-5<-3$.
\item Στους ετερόσημους δεν προσθέτω τις απόλυτες τιμές: $(+9)+(-4)\neq13$, αλλά $9-4=5$.
\item Το πρόσημο του αθροίσματος ετεροσήμων είναι του μεγαλύτερου κατ' απόλυτη τιμή: $(-9)+4=-5$, όχι $+5$.
\item Στο $5+(-3)$ οι παρενθέσεις δείχνουν ότι το $-3$ είναι ο αριθμός που προσθέτω, όχι δύο πράξεις.
\end{itemize}
\end{prosoxi}

\tmima{Λυμένα παραδείγματα}

\paradeigma{Πρόσθεση στον άξονα}
\tixero Οι προσθέσεις $(-2)+5$, $3+(-7)$, $(-4)+(-3)$.\par
\tithelo Να τις κάνω στον άξονα.\par
\possindeo Ξεκινάω από τον πρώτο προσθετέο και κινούμαι: από το $-2$ πέντε δεξιά: $3$. Από το $3$ επτά αριστερά: $-4$. Από το $-4$ τρία αριστερά: $-7$.\par
\apantisi $3$, $-4$, $-7$.\par
\elegxos Με τους κανόνες: $5-2=3$, $-(7-3)=-4$, $-(4+3)=-7$.

\paradeigma{Με τους κανόνες}
\tixero Οι αριθμοί $(-8)+(-5)$, $(+9)+(-4)$, $(-9)+(+4)$, $\left(-\frac{1}{2}\right)+\left(-\frac{1}{4}\right)$.\par
\tithelo Να υπολογίσω τα αθροίσματα.\par
\possindeo Ομόσημοι: $-(8+5)=-13$. Ετερόσημοι με μεγαλύτερη απόλυτη τιμή το $+9$: $+(9-4)=5$. Ετερόσημοι με μεγαλύτερη το $-9$: $-(9-4)=-5$. Κλάσματα ομόσημα: $-\left(\frac{2}{4}+\frac{1}{4}\right)=-\frac{3}{4}$.\par
\apantisi $-13$, $5$, $-5$, $-\frac{3}{4}$.\par
\elegxos Κάθε πρόσημο το βρίσκω από τον προσθετέο με τη μεγαλύτερη απόλυτη τιμή.

\paradeigma{Πολλοί προσθετέοι}
\tixero Το άθροισμα $(-7)+12+(-5)+9+(-3)$.\par
\tithelo Να το υπολογίσω.\par
\possindeo Ομαδοποιώ: θετικοί $12+9=21$. Αρνητικοί $(-7)+(-5)+(-3)=-15$. Τότε $21+(-15)=6$.\par
\apantisi $6$.\par
\elegxos Με τη σειρά: $-7+12=5$, $5-5=0$, $0+9=9$, $9-3=6$.

\paradeigma{Οι ιδιότητες βοηθούν}
\tixero Το άθροισμα $(-4{,}6)+(+3{,}9)+(+4{,}6)+(-3{,}9)$.\par
\tithelo Να το βρω εύκολα.\par
\possindeo Με την αντιμεταθετική και την προσεταιριστική ομαδοποιώ τους αντίθετους: $\big((-4{,}6)+(+4{,}6)\big)+\big((+3{,}9)+(-3{,}9)\big)=0+0=0$.\par
\apantisi $0$.\par
\elegxos Ο καθένας έχει τον αντίθετό του και ακυρώνεται.

\paradeigma{Το ταμείο της παρέας}
\tixero Το ταμείο έχει $+30$ €. Κινήσεις: $-12$, $+8$, $-35$, $+20$.\par
\tithelo Το υπόλοιπο.\par
\possindeo Θετικοί: $30+8+20=58$. Αρνητικοί: $-12-35=-47$. $58+(-47)=11$.\par
\apantisi $11$ €.\par
\elegxos Βήμα βήμα: $30,\ 18,\ 26,\ -9,\ 11$. Συμφωνεί.

\tmima{Ασκήσεις}

\zoni{Το κατάλαβα}{κανόνες, ιδιότητες, πολλοί προσθετέοι}

\ex Στον άξονα βρες τα αθροίσματα: α) $(-2)+5$ \quad β) $3+(-7)$ \quad γ) $(-4)+(-3)$ \quad δ) $6+(-6)$ \quad ε) $(-1)+(+1)$.

\ex Υπολόγισε: α) $(-8)+(-5)$ \quad β) $(+9)+(-4)$ \quad γ) $(-9)+(+4)$ \quad δ) $(+7)+(+5)$ \quad ε) $(-12)+0$ \quad στ) $(-2{,}5)+(+1{,}5)$.

\ex Υπολόγισε: α) $(-3{,}5)+(+3{,}5)$ \quad β) $(-0{,}7)+(-1{,}8)$ \quad γ) $(+2{,}4)+(-5{,}6)$ \quad δ) $\left(-\frac{1}{2}\right)+\left(-\frac{1}{4}\right)$ \quad ε) $\left(+\frac{2}{3}\right)+\left(-\frac{1}{2}\right)$ \quad στ) $\left(-\frac{5}{6}\right)+\left(+\frac{1}{3}\right)$.

\ex Συμπλήρωσε τον αριθμό που λείπει ώστε το άθροισμα να είναι $0$: α) $(-7)+\kouti$ \quad β) $(+3{,}5)+\kouti$ \quad γ) $\kouti+(-4)$ \quad δ) $\left(-\frac{1}{3}\right)+\kouti$.

\ex Υπολόγισε: α) $(-7)+12+(-5)+9+(-3)$ \quad β) $(-1{,}5)+4+(-2{,}5)+1$ \quad γ) $10+(-3)+(-4)+(-5)$.

\ex Υπολόγισε με τις ιδιότητες: α) $(-4{,}6)+(+3{,}9)+(+4{,}6)+(-3{,}9)$ \quad β) $37+(-18)+(-37)+8$.

\ex Σωστό ή λάθος; α) Το άθροισμα δύο αρνητικών είναι αρνητικό. \quad β) Το άθροισμα δύο ετεροσήμων είναι πάντα αρνητικό. \quad γ) $(-5)+5=0$. \quad δ) $(-3)+(-4)=+7$. \quad ε) $\alpha+0=\alpha$. \quad στ) Το άθροισμα δύο αντίθετων είναι $0$. \quad ζ) $(-2)+(+5)>0$.

\zoni{Το χρησιμοποιώ}{θερμόμετρα, ταμεία, παιχνίδια}

\ex \textbf{Η θερμοκρασία.} Το πρωί στο Μέτσοβο είναι $-6$°, το μεσημέρι ανεβαίνει $9$ βαθμούς και το βράδυ πέφτει $5$ βαθμούς. Πόση είναι η θερμοκρασία το μεσημέρι και το βράδυ;

\ex \textbf{Το ταμείο.} Το ταμείο έχει $+30$ €. Γίνονται οι κινήσεις $-12$, $+8$, $-35$, $+20$. Πόσο είναι το υπόλοιπο μετά από κάθε κίνηση και στο τέλος;

\ex \textbf{Ο ανελκυστήρας.} Ο Νίκος είναι στον όροφο $-2$, ανεβαίνει $5$ ορόφους και μετά κατεβαίνει $4$. Σε ποιον όροφο φτάνει;

\ex \textbf{Ο δύτης.} Ο δύτης είναι στα $-18$ μ., ανεβαίνει $7$ μ. και μετά κατεβαίνει $4$ μ. Πού βρίσκεται;

\ex \textbf{Το παιχνίδι.} Οι πόντοι της Πολυξένης στους τέσσερις γύρους είναι $+8$, $-5$, $-7$, $+12$. Πόσους έχει συνολικά;

\ex \textbf{Ο λογαριασμός.} Η Μυρτώ έχει υπόλοιπο $-45$ € στον λογαριασμό, καταθέτει $80$ € και μετά αναλαμβάνει $20$ €. Πόσα έχει;

\ex \textbf{Η διαφορά καλαθιών.} Η ομάδα μπάσκετ έχει στους πέντε αγώνες διαφορά $+12$, $-7$, $-9$, $+4$, $+6$. Ποια είναι η συνολική διαφορά;

\zoni{Το σκέφτομαι}{γιατί ισχύουν οι κανόνες}

\ex Πότε το άθροισμα δύο ρητών είναι $0$;

\ex Αν οι $\alpha,\,\beta$ είναι ετερόσημοι, ποιο πρόσημο έχει το άθροισμά τους;

\ex Γράψε δύο διαφορετικά ζεύγη ρητών με άθροισμα $-5$.

\ex Είναι πάντα το άθροισμα μεγαλύτερο από κάθε προσθετέο; Δώσε αντιπαράδειγμα.

\ex Δείξε με τα $\alpha=3$ και $\beta=4$ ότι $(-\alpha)+(-\beta)=-(\alpha+\beta)$.

\zoni{Για το σπίτι}{ο προϋπολογισμός σου}

\ex Κράτησε για μια εβδομάδα τα έσοδα και έξοδα του χαρτζιλικιού σου ως ρητούς και υπολόγισε το τελικό υπόλοιπο.

\ex Φτιάξε ένα παιχνίδι με κάρτες $+$ και $-$ και παίξε με έναν φίλο τρεις γύρους. Ποιος έχει το μεγαλύτερο άθροισμα;

\paragrafos{Αφαίρεση ρητών: η απόσταση ανάμεσα σε δύο θέσεις}{Α.7.4, «Αφαίρεση ρητών αριθμών»}

\begin{anagke}
Ο μετεωρολόγος στην τηλεόραση του ξενοδοχείου λέει: «Η μέγιστη θερμοκρασία σήμερα είναι $+5$ βαθμοί και η ελάχιστη $-3$. Η διαφορά είναι $8$ βαθμοί». Ο Στέλιος κάνει τον λογαριασμό: «Πέντε μείον τρία κάνει δύο. Ο μετεωρολόγος κάνει λάθος». Η Πολυξένη τον διορθώνει: «Δεν αφαιρούμε το τρία από το πέντε. Αφαιρούμε το μείον τρία!».

Η Αγγελική σχεδιάζει τον άξονα: «Από το $-3$ ως το $+5$ πόσα βήματα υπάρχουν;». Μετράνε: τρία βήματα ως το μηδέν και πέντε ακόμη. Οκτώ. Ο Νίκος συμπληρώνει: «Η αφαίρεση είναι αναποδογύρισμα της πρόσθεσης: ψάχνω αριθμό που προστιθέμενος στο $-3$ να δώσει $+5$. Δηλαδή $(-3)+x=5$. Αφού $(-3)+8=5$, το $5-(-3)$ είναι $8$».

Ο Ηλίας το γενικεύει: «Αν η αφαίρεση ψάχνει «πόσο πρέπει να προσθέσω», μήπως $\alpha-\beta$ είναι το ίδιο με το να προσθέσω τον αντίθετο του $\beta$;». Έτσι η παρέα ανακαλύπτει ότι κάθε αφαίρεση μπορεί να γίνει πρόσθεση και ότι όλα όσα ξέρουν για την πρόσθεση λειτουργούν και για την αφαίρεση.
\end{anagke}

\begin{lexi}{αφαίρεση, διαφορά}{αφαίρεση, διαφορά}
Η \lexiword{αφαίρεση} είναι το «αφ-αιρώ», «βγάζω από κάτι». Η \lexiword{διαφορά} είναι αυτό που μένει και δείχνει «πόσο διαφέρουν» δύο αριθμοί. Ο \lexiword{μειωτέος} είναι αυτός που «μειώνεται» και ο \lexiword{αφαιρετέος} αυτός που αφαιρείται.
\end{lexi}

\begin{metafrasi}
Η διαφορά $\alpha-\beta$ είναι ο αριθμός $x$ που... & $\beta+x=\alpha$\\
Αφαίρεση σημαίνει πρόσθεση του αντίθετου. & $\alpha-\beta=\alpha+(-\beta)$\\
Η διαφορά θερμοκρασιών $+5$ και $-3$. & $(+5)-(-3)=8$\\
Αφαιρώ ολόκληρη παρένθεση. & $\alpha-(\beta+\gamma)=\alpha-\beta-\gamma$\\
\end{metafrasi}

\tmima{Η θεωρία}

\begin{orismos}{διαφορά}
\textbf{Διαφορά} $\alpha-\beta$ δύο ρητών ορίζεται ως η μοναδική λύση της εξίσωσης $\beta+x=\alpha$.
\end{orismos}

Για παράδειγμα η διαφορά $(+5)-(-3)$ είναι η λύση της $(-3)+x=5$, δηλαδή $x=8$. Από τον ορισμό προκύπτει ο κανόνας:

\begin{kanonas}{αφαίρεση ρητών}
Για να αφαιρέσω έναν αριθμό, προσθέτω τον \textbf{αντίθετό} του: $$\alpha-\beta=\alpha+(-\beta).$$ Έτσι $7-(+9)=7+(-9)=-2$ και $(-4)-(-6)=(-4)+(+6)=2$.
\end{kanonas}

Η διαφορά $\alpha-\beta$ δείχνει και τον δρόμο από το $\beta$ στο $\alpha$ στον άξονα: αν $\alpha>\beta$ είναι θετική, αλλιώς αρνητική.

\tmimasmall{Αλγεβρικό άθροισμα και απαλοιφή παρενθέσεων}

Μια παράσταση με προσθέσεις και αφαιρέσεις γράφεται ως πρόσθεση ρητών. Για παράδειγμα $12-15+8-20+3=12+(-15)+8+(-20)+3$, και μπορώ να τα αθροίσω με όποια σειρά θέλω: $(12+8+3)+(-15-20)=23-35=-12$. Τέτοιο άθροισμα ονομάζεται \textbf{αλγεβρικό άθροισμα}. Για να σβήσω παρενθέσεις:
\begin{itemize}[leftmargin=1.2em,itemsep=0.3ex,topsep=2pt]
\item Αν μπροστά από την παρένθεση υπάρχει $+$, τη σβήνω χωρίς να αλλάξω τα πρόσημα: $\alpha+(\beta-\gamma)=\alpha+\beta-\gamma$.
\item Αν μπροστά υπάρχει $-$, τη σβήνω και \textbf{αλλάζω το πρόσημο όλων} των όρων μέσα: $\alpha-(\beta-\gamma+\delta)=\alpha-\beta+\gamma-\delta$.
\end{itemize}

\begin{prosoxi}
\begin{itemize}[leftmargin=1.2em,itemsep=0.3ex,topsep=2pt]
\item Η διαφορά δεν είναι αντιμεταθετική: $5-3=2$ αλλά $3-5=-2$. Δεν ισχύει $\alpha-\beta=\beta-\alpha$.
\item Η αφαίρεση δεν «μικραίνει» πάντα: $3-(-2)=5$, μεγαλύτερο από το $3$.
\item Με το $-$ μπροστά από παρένθεση αλλάζω το πρόσημο \emph{όλων} των όρων της παρένθεσης, όχι μόνο του πρώτου: $10-(3-8+2)=10-3+8-2$.
\item Η διαφορά θερμοκρασιών $+5$ και $-3$ είναι $8$, όχι $2$. Αφαιρώ τον αριθμό μαζί με το πρόσημό του.
\end{itemize}
\end{prosoxi}

\tmima{Λυμένα παραδείγματα}

\paradeigma{Η διαφορά των θερμοκρασιών}
\tixero Μέγιστη $+5$°, ελάχιστη $-3$°.\par
\tithelo Η θερμική διακύμανση, δηλαδή η διαφορά μέγιστης και ελάχιστης.\par
\possindeo $(+5)-(-3)=(+5)+(+3)=8$. Έλεγχος στον άξονα: από το $-3$ ως το $+5$ υπάρχουν $8$ μονάδες.\par
\apantisi $8$ βαθμοί.\par
\elegxos $(-3)+8=5$.

\paradeigma{Με τον κανόνα}
\tixero Οι διαφορές $7-(+9)$, $(-4)-(-6)$, $(+3{,}5)-(-1{,}5)$, $\left(-\frac{1}{2}\right)-\left(+\frac{1}{3}\right)$.\par
\tithelo Να τις υπολογίσω.\par
\possindeo Προσθέτω τον αντίθετο: $7+(-9)=-2$. $(-4)+(+6)=2$. $3{,}5+1{,}5=5$. $-\frac{1}{2}-\frac{1}{3}=-\frac{3}{6}-\frac{2}{6}=-\frac{5}{6}$.\par
\apantisi $-2$, $2$, $5$, $-\frac{5}{6}$.\par
\elegxos Για την πρώτη: $9+(-2)=7$. Σωστό.

\paradeigma{Απαλοιφή παρενθέσεων}
\tixero Η παράσταση $10-(3-8+2)$.\par
\tithelo Να την υπολογίσω με δύο τρόπους.\par
\possindeo \emph{1ος τρόπος:} πρώτα η παρένθεση: $3-8+2=-3$, και $10-(-3)=13$. \emph{2ος τρόπος:} σβήνω την παρένθεση αλλάζοντας πρόσημα: $10-3+8-2=13$.\par
\apantisi $13$.\par
\elegxos Και οι δύο δίνουν $13$.

\paradeigma{Αλγεβρικό άθροισμα}
\tixero Η παράσταση $12-15+8-20+3$.\par
\tithelo Να την υπολογίσω.\par
\possindeo Ομαδοποιώ: θετικοί $12+8+3=23$. Αρνητικοί $-15-20=-35$. $23+(-35)=-12$.\par
\apantisi $-12$.\par
\elegxos Με τη σειρά: $12-15=-3$, $-3+8=5$, $5-20=-15$, $-15+3=-12$.

\paradeigma{Εξισώσεις με αφαίρεση}
\tixero Οι εξισώσεις $x+(-3)=5$, $x-7=-2$ και $4-x=9$.\par
\tithelo Να τις λύσω.\par
\possindeo $x=5-(-3)=8$. $x=-2+7=5$. $x=4-9=-5$ (από την $\alpha-x=\beta\Rightarrow x=\alpha-\beta$).\par
\apantisi $8$, $5$, $-5$.\par
\elegxos $8+(-3)=5$, $5-7=-2$, $4-(-5)=9$.

\tmima{Ασκήσεις}

\zoni{Το κατάλαβα}{διαφορές, παρενθέσεις, αλγεβρικά αθροίσματα}

\ex Υπολόγισε: α) $(+5)-(+3)$ \quad β) $(+3)-(+5)$ \quad γ) $(-4)-(-1)$ \quad δ) $(-1)-(-4)$ \quad ε) $0-(+6)$ \quad στ) $0-(-6)$.

\ex Μετάτρεψε σε πρόσθεση και υπολόγισε: α) $7-(+9)$ \quad β) $(-4)-(-6)$ \quad γ) $(+3{,}5)-(-1{,}5)$ \quad δ) $\left(-\frac{1}{2}\right)-\left(+\frac{1}{3}\right)$ \quad ε) $(-2{,}4)-(+0{,}6)$ \quad στ) $8-(-8)$.

\ex Λύσε: α) $x+7=3$ \quad β) $x-5=-2$ \quad γ) $4-x=9$ \quad δ) $x+(-6)=-6$ \quad ε) $-3-x=2$.

\ex Κάνε απαλοιφή παρενθέσεων και υπολόγισε: α) $10-(3-8+2)$ \quad β) $5+(-3+7)-(2-9)$ \quad γ) $-(-4+5)-(7-10)$ \quad δ) $20-(-5)-(-(3))$.

\ex Υπολόγισε τα αλγεβρικά αθροίσματα: α) $12-15+8-20+3$ \quad β) $-2{,}5+7-4{,}5-1$ \quad γ) $\frac{1}{2}-\frac{3}{4}+\frac{1}{4}-1$.

\ex Σωστό ή λάθος; α) $\alpha-\beta=\beta-\alpha$. \quad β) $\alpha-0=\alpha$. \quad γ) $0-\alpha=-\alpha$. \quad δ) $(-3)-(-3)=0$. \quad ε) $(+4)-(-2)=2$. \quad στ) Η διαφορά δύο θετικών μπορεί να είναι αρνητική.

\zoni{Το χρησιμοποιώ}{θερμοκρασίες, υψόμετρα, χρόνια}

\ex \textbf{Η διακύμανση.} Βρες τη διαφορά μέγιστης και ελάχιστης θερμοκρασίας όταν α) μέγιστη $+5$, ελάχιστη $-3$ \quad β) μέγιστη $-2$, ελάχιστη $-11$.

\ex \textbf{Το υψόμετρο.} Ο Όλυμπος είναι στα $+2\,917$ μ. και η Νεκρά Θάλασσα στα $-430$ μ. Πόση είναι η υψομετρική διαφορά;

\ex \textbf{Η ιστορία.} Η δυναστεία των Χαν ξεκίνησε το $-206$ και ο Διόφαντος έζησε γύρω στο $250$ μ.Χ. Πόσα χρόνια μεσολαβούν (αγνοώντας ότι δεν υπήρξε έτος $0$);

\ex \textbf{Το ταμείο.} Η Ζωή έχει $-15$ € και ο Ηλίας $+8$ €. Πόση είναι η διαφορά των δύο ποσών;

\ex \textbf{Ο δύτης.} Ο δύτης ανεβαίνει από το $-22$ μ. στο $-9$ μ. Πόσα μέτρα ανέβηκε;

\ex \textbf{Ο ανελκυστήρας.} Από τον όροφο $3$ στον $-2$ η μεταβολή είναι $(-2)-3$. Πόσο είναι και τι σημαίνει;

\ex \textbf{Η πτώση.} Η θερμοκρασία από $+4$° έπεσε στα $-7$°. Πόση είναι η μεταβολή (νέα $-$ παλιά);

\zoni{Το σκέφτομαι}{όταν η αφαίρεση μεγαλώνει}

\ex Ο Στέλιος γράφει $5-(-3)=2$. Βρες το λάθος και δείξε το σωστό με τον άξονα.

\ex Μικραίνει πάντα ένας αριθμός όταν του αφαιρώ κάτι; Δώσε παράδειγμα.

\ex Αν $\alpha>\beta$, ποιο πρόσημο έχει η διαφορά $\alpha-\beta$; Και αν $\alpha<\beta$;

\ex Δείξε με $\alpha=10$, $\beta=3$, $\gamma=4$ ότι $\alpha-(\beta+\gamma)=\alpha-\beta-\gamma$.

\ex Βρες τον αριθμό που, αν του αφαιρέσουμε το $-5$, δίνει $2$.

\zoni{Για το σπίτι}{διαφορές του κόσμου}

\ex Βρες τη διαφορά θερμοκρασίας ανάμεσα στο πιο ζεστό και στο πιο κρύο μέρος του κόσμου σήμερα.

\ex Φτιάξε μια παράσταση με έξι όρους και παρενθέσεις και υπολόγισέ την με δύο τρόπους.

\paragrafos{Πολλαπλασιασμός ρητών: το μείον επί μείον}{Α.7.5, «Πολλαπλασιασμός ρητών αριθμών»}

\begin{anagke}
Το βράδυ στο Μέτσοβο η θερμοκρασία πέφτει $2$ βαθμούς κάθε ώρα. Η Πολυξένη παρατηρεί το θερμόμετρο στις $20{:}00$, που δείχνει $0$°. «Τι θα δείχνει σε τρεις ώρες;». Ο Νίκος σκέφτεται: «Κάθε ώρα $-2$, τρεις ώρες: $-2-2-2=-6$». Η Αγγελική το γράφει ως πολλαπλασιασμό: $3\cdot(-2)=-6$. «Και πριν από τρεις ώρες, στις $17{:}00$;» ρωτάει η Μυρτώ. «Αφού τώρα πέφτει δύο βαθμούς ανά ώρα, πριν από τρεις ώρες είχε $6$ βαθμούς περισσότερο», απαντάει η Πολυξένη. Ο ρυθμός είναι $-2$ και ο χρόνος «τρεις ώρες πριν» είναι $-3$, και το αποτέλεσμα είναι $+6$. Άρα $(-2)\cdot(-3)=+6$.

Ο Στέλιος ξαφνιάζεται: «Το μείον επί μείον κάνει συν; Σαν τη διπλή άρνηση της γλώσσας που καταλήγει σε κατάφαση;». Η κυρία Αλεξάνδρα του απαντά: «Η γλώσσα είναι ωραία αναλογία, αλλά τα μαθηματικά θέλουν να δουν γιατί». Το ίδιο το μοτίβο δείχνει το γιατί: $3\cdot(-2)=-6$, $2\cdot(-2)=-4$, $1\cdot(-2)=-2$, $0\cdot(-2)=0$. Κάθε φορά που το πρώτο μικραίνει κατά $1$ το γινόμενο μεγαλώνει κατά $2$. Συνεχίζοντας: $(-1)\cdot(-2)=2$, $(-2)\cdot(-2)=4$.
\end{anagke}

\begin{lexi}{παράγοντας, γινόμενο, ομόσημοι}{παράγοντας, γινόμενο}
\lexiword{Παράγοντες} είναι οι αριθμοί που πολλαπλασιάζονται και \lexiword{γινόμενο} αυτό που «γίνεται». Η \lexiword{επιμεριστική} ιδιότητα είναι αυτή που «επιμερίζει» (μοιράζει σε μέρη) τον πολλαπλασιασμό πάνω σε ένα άθροισμα.
\end{lexi}

\begin{metafrasi}
Πέφτει $2$ βαθμούς την ώρα, για $3$ ώρες. & $3\cdot(-2)=-6$\\
Πριν από $3$ ώρες. & $(-2)\cdot(-3)=+6$\\
Ετερόσημοι. & γινόμενο αρνητικό\\
Ομόσημοι. & γινόμενο θετικό\\
Ο πολλαπλασιασμός πάνω σε άθροισμα. & $\alpha(\beta+\gamma)=\alpha\beta+\alpha\gamma$\\
\end{metafrasi}

\tmima{Η θεωρία}

\begin{kanonas}{γινόμενο ρητών}
Το γινόμενο δύο ρητών έχει απόλυτη τιμή το γινόμενο των απόλυτων τιμών τους και:
\begin{itemize}[leftmargin=1.2em,itemsep=0.3ex,topsep=2pt]
\item πρόσημο \textbf{θετικό} αν οι παράγοντες είναι ομόσημοι: $(-7)\cdot(-8)=+56$, $(+5)(+4)=+20$.
\item πρόσημο \textbf{αρνητικό} αν οι παράγοντες είναι ετερόσημοι: $(-6)\cdot(+4)=-24$.
\item Αν ένας παράγοντας είναι $0$, το γινόμενο είναι $0$.
\end{itemize}
\end{kanonas}

Για παράδειγμα $\left(-\frac{3}{4}\right)\cdot\left(-\frac{8}{9}\right)=+\frac{3\cdot8}{4\cdot9}=\frac{2}{3}$ και $(+0{,}5)\cdot(-12)=-6$.

\begin{idiotita}{ιδιότητες του πολλαπλασιασμού}
Για κάθε τρεις ρητούς ισχύουν: \textbf{αντιμεταθετική} $\alpha\beta=\beta\alpha$, \textbf{προσεταιριστική} $(\alpha\beta)\gamma=\alpha(\beta\gamma)$, $\alpha\cdot1=\alpha$, $\alpha\cdot0=0$, $\alpha\cdot(-1)=-\alpha$ και \textbf{επιμεριστική} $\alpha(\beta+\gamma)=\alpha\beta+\alpha\gamma$ (και $\alpha(\beta-\gamma)=\alpha\beta-\alpha\gamma$).
\end{idiotita}

\tmimasmall{Γινόμενο πολλών παραγόντων}

Το γινόμενο πολλών ρητών (όλων μη μηδενικών) είναι θετικό αν το πλήθος των αρνητικών παραγόντων είναι \emph{άρτιο} και αρνητικό αν είναι \emph{περιττό}. Η απόλυτη τιμή του είναι το γινόμενο των απόλυτων τιμών. Έτσι $(-2)\cdot3\cdot(-5)\cdot(-1)=-30$ (τρεις αρνητικοί).

Η επιμεριστική ιδιότητα δίνει συντομίες: $99\cdot(-6)=(100-1)\cdot(-6)=-600+6=-594$.

\begin{prosoxi}
\begin{itemize}[leftmargin=1.2em,itemsep=0.3ex,topsep=2pt]
\item Στον πολλαπλασιασμό, δύο αρνητικοί δίνουν θετικό, ενώ στην πρόσθεση δύο αρνητικοί δίνουν αρνητικό: $(-3)\cdot(-4)=+12$ και $(-3)+(-4)=-7$.
\item Το πρόσημο του γινομένου μετράει το πλήθος των αρνητικών παραγόντων, όχι το πλήθος των θετικών.
\item Στην επιμεριστική πολλαπλασιάζω τον αριθμό με \emph{κάθε} όρο της παρένθεσης: $(-2)(5-9)=-10+18$.
\item $(-3)^2$ και $-3^2$ δεν είναι ίδια, αλλά θα το δούμε στις δυνάμεις.
\end{itemize}
\end{prosoxi}

\tmima{Λυμένα παραδείγματα}

\paradeigma{Το μοτίβο}
\tixero Τα γινόμενα $3\cdot(-4)$, $2\cdot(-4)$, $1\cdot(-4)$, $0\cdot(-4)$.\par
\tithelo Να συνεχίσω το μοτίβο.\par
\possindeo Τα γινόμενα είναι $-12$, $-8$, $-4$, $0$ και αυξάνουν κατά $4$ κάθε φορά που ο πρώτος παράγοντας μειώνεται κατά $1$. Άρα $(-1)\cdot(-4)=4$, $(-2)\cdot(-4)=8$, $(-3)\cdot(-4)=12$.\par
\apantisi $4$, $8$, $12$.\par
\elegxos Ο κανόνας ομόσημοι δίνει θετικό: $(-3)\cdot(-4)=+12$.

\paradeigma{Γινόμενα με τον κανόνα}
\tixero Τα $(-6)(+4)$, $(-7)(-8)$, $(+0{,}5)(-12)$, $\left(-\frac{3}{4}\right)\left(-\frac{8}{9}\right)$.\par
\tithelo Να τα υπολογίσω.\par
\possindeo Πρώτα το πρόσημο, μετά οι απόλυτες τιμές. Ετερόσημοι: $-24$. Ομόσημοι: $+56$. Ετερόσημοι: $-6$. Ομόσημοι: $+\frac{24}{36}=\frac{2}{3}$.\par
\apantisi $-24$, $56$, $-6$, $\frac{2}{3}$.\par
\elegxos Το $7\cdot8=56$ και $6\cdot4=24$.

\paradeigma{Πολλοί παράγοντες}
\tixero Το γινόμενο $(-2)\cdot3\cdot(-5)\cdot(-1)$.\par
\tithelo Να βρω το πρόσημο και την τιμή.\par
\possindeo Αρνητικοί παράγοντες: $3$ (περιττό πλήθος), άρα αρνητικό. Απόλυτη τιμή: $2\cdot3\cdot5\cdot1=30$.\par
\apantisi $-30$.\par
\elegxos Βήμα βήμα: $(-2)\cdot3=-6$, $-6\cdot(-5)=30$, $30\cdot(-1)=-30$.

\paradeigma{Η επιμεριστική}
\tixero Οι $(-4)\cdot(7+3)$ και $99\cdot(-6)$.\par
\tithelo Να τις υπολογίσω εύκολα.\par
\possindeo $(-4)\cdot7+(-4)\cdot3=-28-12=-40$ (ή $-4\cdot10$). Για το δεύτερο: $99=100-1$, $(100-1)\cdot(-6)=-600+6=-594$.\par
\apantisi $-40$ και $-594$.\par
\elegxos $99\cdot6=594$ και ο αριθμός είναι αρνητικός.

\paradeigma{Παράσταση}
\tixero Η παράσταση $3\cdot(-4)-(-5)\cdot2+(-2)\cdot(-3)$.\par
\tithelo Να την υπολογίσω.\par
\possindeo Πρώτα οι πολλαπλασιασμοί: $3\cdot(-4)=-12$, $(-5)\cdot2=-10$, $(-2)(-3)=6$. Άρα $-12-(-10)+6=-12+10+6=4$.\par
\apantisi $4$.\par
\elegxos Προτεραιότητα: πολλαπλασιασμοί πριν από προσθέσεις και αφαιρέσεις.

\tmima{Ασκήσεις}

\zoni{Το κατάλαβα}{πρόσημα και τιμές}

\ex Συμπλήρωσε το μοτίβο: $3\cdot(-4)$, $2\cdot(-4)$, $1\cdot(-4)$, $0\cdot(-4)$, $(-1)\cdot(-4)$, $(-2)\cdot(-4)$, $(-3)\cdot(-4)$.

\ex Υπολόγισε: α) $(-6)\cdot(+4)$ \quad β) $(+7)\cdot(-3)$ \quad γ) $(-5)\cdot(-9)$ \quad δ) $(-12)\cdot0$ \quad ε) $(+8)\cdot(+5)$ \quad στ) $(-1)\cdot(-1)$.

\ex Υπολόγισε: α) $(+0{,}5)\cdot(-12)$ \quad β) $(-1{,}2)\cdot(-3)$ \quad γ) $(-2{,}5)\cdot(+0{,}4)$ \quad δ) $\left(-\frac{3}{4}\right)\cdot\left(-\frac{8}{9}\right)$ \quad ε) $\left(+\frac{2}{5}\right)\cdot(-15)$ \quad στ) $\left(-\frac{1}{2}\right)\cdot\left(-\frac{2}{3}\right)\cdot(-3)$.

\ex Βρες το πρόσημο του γινομένου χωρίς να το υπολογίσεις: α) $(-2)\cdot3\cdot(-5)\cdot(-1)$ \quad β) $(-1)(-1)(-1)(-1)$ \quad γ) $7\cdot(-3)\cdot(-2)\cdot4\cdot(-1)$ \quad δ) $(-5)\cdot0\cdot(-8)$.

\ex Υπολόγισε: α) $(-8)\cdot1$ \quad β) $(-8)\cdot(-1)$ \quad γ) $0\cdot(-3)$ \quad δ) $\alpha\cdot(-1)$ για $\alpha=-7$.

\ex Με την επιμεριστική: α) $(-4)\cdot(7+3)$ \quad β) $5\cdot(-3+8)$ \quad γ) $(-2)\cdot(5-9)$ \quad δ) $99\cdot(-6)$ \quad ε) $(-12)\cdot\left(\frac{1}{3}+\frac{1}{4}\right)$.

\ex Υπολόγισε: α) $3\cdot(-4)-(-5)\cdot2+(-2)\cdot(-3)$ \quad β) $(-2)\cdot(4-7)+3\cdot(-1-2)$ \quad γ) $10-2\cdot(-3)\cdot(-5)$.

\ex Σωστό ή λάθος; α) $(-3)(-4)=-12$. \quad β) Το γινόμενο ετεροσήμων είναι αρνητικό. \quad γ) $\alpha\cdot0=0$. \quad δ) $(-1)\cdot\alpha=-\alpha$. \quad ε) Το γινόμενο τριών αρνητικών είναι θετικό. \quad στ) $(-2)(-3)(-4)=-24$.

\zoni{Το χρησιμοποιώ}{θερμοκρασίες, χρέη, βουτιές}

\ex \textbf{Η θερμοκρασία.} Πέφτει $2$ βαθμούς κάθε ώρα. Πόσο θα έχει μεταβληθεί σε $3$ ώρες και πόσο έχει μεταβληθεί από πριν $3$ ώρες;

\ex \textbf{Το χρέος.} Κάθε εβδομάδα το χρέος της Μυρτώς αυξάνεται κατά $4$ €, δηλαδή η μεταβολή είναι $-4$ €. Ποια είναι η μεταβολή σε $5$ εβδομάδες;

\ex \textbf{Ο δύτης.} Ο δύτης κατεβαίνει $3$ μέτρα το λεπτό. Πού θα είναι σε $8$ λεπτά; Πού ήταν πριν από $5$ λεπτά;

\ex \textbf{Ο ανελκυστήρας.} Ο ανελκυστήρας κατεβαίνει $2$ ορόφους κάθε φορά, $4$ φορές. Πόση είναι η συνολική μεταβολή;

\ex \textbf{Το ταμείο.} Αν τα $9$ παιδιά της παρέας χρωστάνε από $4$ € το καθένα, τι αριθμός εκφράζει το συνολικό χρέος;

\ex \textbf{Η μετοχή.} Η τιμή μιας μετοχής πέφτει $0{,}25$ € κάθε ημέρα. Πόση είναι η συνολική μεταβολή σε $6$ ημέρες;

\zoni{Το σκέφτομαι}{γιατί το μείον επί μείον κάνει συν}

\ex Εξήγησε με το μοτίβο ή με τη θερμοκρασία γιατί $(-2)\cdot(-3)=+6$.

\ex Γιατί το τετράγωνο ενός ρητού δεν είναι ποτέ αρνητικό;

\ex Πόσα αρνητικά πρέπει να έχει ένα γινόμενο $10$ παραγόντων για να είναι θετικό; Δώσε δύο διαφορετικά πλήθη.

\ex Αν $\alpha\cdot\beta=0$, τι συμπεραίνεις για τους $\alpha,\,\beta$;

\ex Υπολόγισε το $(-1)\cdot(-1)\cdot\ldots\cdot(-1)$ με $100$ παράγοντες και με $101$ παράγοντες.

\zoni{Για το σπίτι}{φτιάξε το μοτίβο}

\ex Φτιάξε το μοτίβο του πολλαπλασιασμού με το $-5$ από $3$ ως $-3$ και πες τι παρατηρείς.

\paragrafos{Διαίρεση ρητών: μοιράζω και τα χρέη}{Α.7.6, «Διαίρεση ρητών αριθμών»}

\begin{anagke}
Στο τέλος της εκδρομής ο Νίκος, που κρατάει το ταμείο, ανακοινώνει με σκυμμένο κεφάλι: «Το ταμείο έμεινε στο $-24$ ευρώ. Χρωστάμε $24$ ευρώ». Η Μυρτώ ρωτάει: «Αν το μοιράσουμε εξίσου στα έξι παιδιά της παρέας, πόσο ήταν το καθένα;». Η Πολυξένη το γράφει: $(-24):6$. Ο Βασίλης λέει «τέσσερα» και ο Στέλιος διαμαρτύρεται: «Πόσο, χρέος ή κέρδος;». Αφού το χρέος μοιράζεται, κάθε παιδί χρωστά $4$ ευρώ, δηλαδή το πηλίκο είναι $-4$.

Ο Ηλίας ρωτάει το αντίστροφο: «Αν κάθε παιδί πρέπει να βάλει $3$ ευρώ ($-3$ στο ταμείο του) και το σύνολο είναι $-24$, πόσα παιδιά είναι;». Ο Νίκος σκέφτεται: «Πόσες φορές χωράει το $-3$ στο $-24$; Οκτώ φορές». Άρα $(-24):(-3)=+8$. Η Αγγελική τα βάζει σε εξίσωση: $(-3)\cdot x=-24$ και $x=8$. Και η Πολυξένη θυμάται ότι η διαίρεση είναι ο αντίστροφος δρόμος του πολλαπλασιασμού, όπως η αφαίρεση της πρόσθεσης.

Από τον πολλαπλασιασμό ήδη ξέρουν τους κανόνες των προσήμων. Τους ίδιους κανόνες θα βρουν και για τη διαίρεση, και θα δουν γιατί, πάλι, δεν μπορούν να διαιρέσουν με το $0$.
\end{anagke}

\begin{lexi}{διαίρεση, πηλίκο, αντίστροφος}{διαίρεση, πηλίκο, λόγος}
Η \lexiword{διαίρεση} είναι το «δι-αιρώ», δηλαδή «χωρίζω σε μέρη». Ο \lexiword{διαιρετέος} είναι αυτός που μοιράζεται, ο \lexiword{διαιρέτης} αυτός που μοιράζει και το \lexiword{πηλίκο} (από το «πηλίκος», «πόσος») το αποτέλεσμα. Ο \lexiword{λόγος} δύο ρητών είναι το πηλίκο τους, γραμμένο ως κλάσμα.
\end{lexi}

\begin{metafrasi}
Η διαίρεση $\alpha:\beta$ είναι η λύση της εξίσωσης. & $\beta\cdot x=\alpha$ (με $\beta\neq0$)\\
Μοιράζω χρέος $24$ € σε $6$ παιδιά. & $(-24):6=-4$\\
Πόσες φορές χωράει το $-3$ στο $-24$; & $(-24):(-3)=+8$\\
Διαίρεση σημαίνει πολλαπλασιασμός με τον αντίστροφο. & $\alpha:\beta=\alpha\cdot\dfrac{1}{\beta}$\\
Κλάσμα με αρνητικό. & $\dfrac{-6}{3}=\dfrac{6}{-3}=-2$\\
\end{metafrasi}

\tmima{Η θεωρία}

\begin{orismos}{πηλίκο}
\textbf{Πηλίκο} $\alpha:\beta$ δύο ρητών με $\beta\neq0$ ορίζεται ως η μοναδική λύση της εξίσωσης $\beta\cdot x=\alpha$.
\end{orismos}

Για παράδειγμα το $12:(-3)$ είναι η λύση της $(-3)\cdot x=12$, δηλαδή $x=-4$. Από αυτό προκύπτει ότι τα πρόσημα είναι τα ίδια με του πολλαπλασιασμού:

\begin{kanonas}{διαίρεση ρητών}
Το πηλίκο δύο ρητών έχει απόλυτη τιμή το πηλίκο των απόλυτων τιμών και \textbf{πρόσημο} θετικό αν είναι ομόσημοι και αρνητικό αν είναι ετερόσημοι. Δηλαδή $\alpha:\beta$ και $\alpha\cdot\beta$ είναι \textbf{ομόσημοι}. Το $0:\beta=0$. Δεν διαιρούμε με το $0$.
\end{kanonas}

Έτσι $(-24):6=-4$, $(-24):(-3)=+8$, $18:(-6)=-3$, $0:(-5)=0$.

\tmimasmall{Διαίρεση και αντίστροφος}

Ο \textbf{αντίστροφος} ενός ρητού $\alpha\neq0$ είναι ο $\frac{1}{\alpha}$ (γινόμενο $1$). Ο αντίστροφος του $-4$ είναι ο $-\frac{1}{4}$, του $-\frac{2}{3}$ ο $-\frac{3}{2}$. Το $0$ δεν έχει αντίστροφο. Κάθε διαίρεση γίνεται πολλαπλασιασμός: $\alpha:\beta=\alpha\cdot\frac{1}{\beta}$, όπως και στα κλάσματα. Έτσι $\left(-\frac{3}{4}\right):\frac{1}{2}=\left(-\frac{3}{4}\right)\cdot2=-\frac{3}{2}$.

\tmimasmall{Το πηλίκο ως λόγος}

Το πηλίκο δύο ρητών γράφεται ως κλάσμα (\emph{λόγος}): $\dfrac{-6}{3}=-2$, $\dfrac{6}{-3}=-2$, $\dfrac{-6}{-3}=+2$. Ισχύει $\dfrac{-\alpha}{\beta}=\dfrac{\alpha}{-\beta}=-\dfrac{\alpha}{\beta}$. Τα σύνθετα κλάσματα λύνονται όπως πριν: $\dfrac{-\frac{6}{5}}{\frac{3}{10}}=-\dfrac{6}{5}\cdot\dfrac{10}{3}=-4$.

\begin{prosoxi}
\begin{itemize}[leftmargin=1.2em,itemsep=0.3ex,topsep=2pt]
\item Δεν υπάρχει $\alpha:0$: η $0\cdot x=\alpha$ δεν έχει λύση όταν $\alpha\neq0$ και έχει άπειρες όταν $\alpha=0$.
\item Το $(-12):4$ είναι $-3$, όχι $3$: ετερόσημοι δίνουν αρνητικό.
\item Η διαίρεση δεν είναι αντιμεταθετική: $(-6):3=-2$ αλλά $3:(-6)=-\frac{1}{2}$.
\item Αντίστροφος δεν είναι αντίθετος: ο αντίθετος του $4$ είναι $-4$ και ο αντίστροφος $\frac{1}{4}$.
\end{itemize}
\end{prosoxi}

\tmima{Λυμένα παραδείγματα}

\paradeigma{Το ταμείο μοιράζεται}
\tixero Το ταμείο είναι $-24$ € και μοιράζεται σε $6$ παιδιά.\par
\tithelo Το μερίδιο του καθενός και πόσα παιδιά βάζουν από $-3$.\par
\possindeo $(-24):6=-4$ (ετερόσημοι, αρνητικό). $(-24):(-3)=+8$ (ομόσημοι, θετικό).\par
\apantisi $-4$ € και $8$ παιδιά.\par
\elegxos $6\cdot(-4)=-24$ και $8\cdot(-3)=-24$.

\paradeigma{Με τον ορισμό}
\tixero Το πηλίκο $12:(-3)$.\par
\tithelo Να το βρω ως λύση εξίσωσης.\par
\possindeo Ψάχνω $x$ με $(-3)\cdot x=12$. Αφού $(-3)\cdot(-4)=12$, είναι $x=-4$.\par
\apantisi $-4$.\par
\elegxos Ο κανόνας: ετερόσημοι, αρνητικό, $12:3=4$.

\paradeigma{Πηλίκα ρητών}
\tixero Τα $(-7{,}5):(-2{,}5)$, $\left(-\frac{3}{4}\right):\frac{1}{2}$, $(+0{,}6):(-1{,}5)$.\par
\tithelo Να τα υπολογίσω.\par
\possindeo Ομόσημοι: $7{,}5:2{,}5=3$, άρα $+3$. Ετερόσημοι: $-\frac{3}{4}\cdot2=-\frac{3}{2}$. Ετερόσημοι: $0{,}6:1{,}5=0{,}4$, άρα $-0{,}4$.\par
\apantisi $3$, $-\frac{3}{2}$, $-0{,}4$.\par
\elegxos $(-2{,}5)\cdot3=-7{,}5$ και $(-1{,}5)\cdot(-0{,}4)=0{,}6$.

\paradeigma{Αντίστροφοι}
\tixero Οι αριθμοί $-4$, $-\frac{2}{3}$, $5$, $0$.\par
\tithelo Οι αντίστροφοί τους.\par
\possindeo Ο αντίστροφος του $\alpha$ είναι ο $\frac{1}{\alpha}$ με το ίδιο πρόσημο. Του $-4$: $-\frac{1}{4}$. Του $-\frac{2}{3}$: $-\frac{3}{2}$. Του $5$: $\frac{1}{5}$. Το $0$ δεν έχει.\par
\apantisi $-\frac{1}{4}$, $-\frac{3}{2}$, $\frac{1}{5}$ και το $0$ δεν έχει αντίστροφο.\par
\elegxos $(-4)\cdot\left(-\frac{1}{4}\right)=1$.

\paradeigma{Παραστάσεις και σύνθετο κλάσμα}
\tixero Οι παραστάσεις $(-12):3+(-5)\cdot(-2)$ και $\dfrac{-\frac{6}{5}}{\frac{3}{10}}$.\par
\tithelo Να τις υπολογίσω.\par
\possindeo $(-12):3=-4$ και $(-5)(-2)=10$, άρα $-4+10=6$. Το σύνθετο: $-\frac{6}{5}\cdot\frac{10}{3}=-\frac{60}{15}=-4$.\par
\apantisi $6$ και $-4$.\par
\elegxos $-4\cdot\frac{3}{10}=-\frac{6}{5}$.

\tmima{Ασκήσεις}

\zoni{Το κατάλαβα}{πηλίκα, εξισώσεις, αντίστροφοι}

\ex Υπολόγισε: α) $(-24):6$ \quad β) $(-24):(-3)$ \quad γ) $18:(-6)$ \quad δ) $0:(-5)$ \quad ε) $(-7{,}5):(-2{,}5)$ \quad στ) $(+36):(+9)$.

\ex Λύσε: α) $(-3)x=12$ \quad β) $5x=-20$ \quad γ) $x:(-2)=6$ \quad δ) $(-4)x=-4$ \quad ε) $x:3=-5$ \quad στ) $12:x=-4$.

\ex Υπολόγισε: α) $\left(-\frac{3}{4}\right):\frac{1}{2}$ \quad β) $\left(-\frac{2}{3}\right):\left(-\frac{4}{9}\right)$ \quad γ) $5:\left(-\frac{2}{5}\right)$ \quad δ) $(-1{,}2):0{,}4$ \quad ε) $(+0{,}6):(-1{,}5)$ \quad στ) $\left(-\frac{1}{2}\right):\left(-\frac{1}{2}\right)$.

\ex Βρες τους αντίστροφους των $-4$, $-\frac{2}{3}$, $5$, $-0{,}5$, $\frac{1}{3}$.

\ex Απλοποίησε τα σύνθετα κλάσματα: α) $\dfrac{-\frac{6}{5}}{\frac{3}{10}}$ \quad β) $\dfrac{\frac{2}{3}}{-\frac{4}{9}}$ \quad γ) $\dfrac{-5}{-\frac{1}{2}}$.

\ex Υπολόγισε: α) $(-12):3+(-5)\cdot(-2)$ \quad β) $20:(-4)-(-3)\cdot2$ \quad γ) $[(-18):(-6)]\cdot(-4)$.

\ex Πες το πρόσημο χωρίς να υπολογίσεις: α) $(-45):9$ \quad β) $(-45):(-9)$ \quad γ) $45:(-9)$ \quad δ) $(-1\,000):(-1)$.

\ex Σωστό ή λάθος; α) $(-12):(-4)=-3$. \quad β) $0:(-5)=0$. \quad γ) $(-5):0=0$. \quad δ) $\alpha:\beta=\alpha\cdot\frac{1}{\beta}$. \quad ε) Το πηλίκο ετεροσήμων είναι αρνητικό. \quad στ) $\frac{-6}{3}=\frac{6}{-3}$.

\zoni{Το χρησιμοποιώ}{χρέη, ρυθμοί, μέσοι όροι}

\ex \textbf{Το ταμείο.} Το ταμείο είναι $-24$ € και μοιράζεται εξίσου στα $6$ παιδιά. Πόσο αντιστοιχεί στο καθένα; Αν κάθε παιδί βάζει $-3$ €, πόσα παιδιά έχουν συνεισφέρει συνολικά $-24$ €;

\ex \textbf{Η θερμοκρασία.} Η θερμοκρασία έπεσε $18$ βαθμούς σε $6$ ώρες. Ποιος είναι ο ρυθμός μεταβολής ανά ώρα; Και αν έπεσε $12$ βαθμούς σε $4$ ώρες;

\ex \textbf{Ο δύτης.} Ο δύτης κατέβηκε από το $0$ στο $-36$ μ. σε $9$ λεπτά. Πόσα μέτρα το λεπτό;

\ex \textbf{Σε πόση ώρα;} Η θερμοκρασία πέφτει $2{,}5$ βαθμούς την ώρα. Σε πόσες ώρες θα πέσει $15$ βαθμούς;

\ex \textbf{Ο μέσος όρος.} Οι θερμοκρασίες μιας εβδομάδας ήταν $-3$, $-5$, $2$, $0$, $-4$, $1$, $-1$. Βρες το άθροισμα και τον μέσο όρο στο δέκατο.

\ex \textbf{Οι πόντοι.} Οι διαφορές πόντων της Φανής σε $4$ αγώνες ήταν $-8$, $+4$, $-10$, $+6$. Ποιος είναι ο μέσος όρος;

\ex \textbf{Τα παιδιά.} Το χρέος των $-45$ € μοιράστηκε εξίσου και κάθε παιδί χρωστά $-5$ €. Πόσα είναι τα παιδιά;

\zoni{Το σκέφτομαι}{γιατί όχι με το μηδέν}

\ex Γιατί δεν διαιρούμε με το $0$; Εξέτασε με τον ορισμό τις $0\cdot x=5$ και $0\cdot x=0$.

\ex Ο Στέλιος γράφει $(-12):4=3$. Βρες το λάθος και έλεγξε το σωστό με πολλαπλασιασμό.

\ex Σύγκρινε τα πρόσημα των $(-6):3$, $6:(-3)$, $(-6):(-3)$.

\ex Αν $\alpha:\beta=2$ και $\alpha=-10$, βρες το $\beta$.

\ex Δείξε με $\alpha=6$ και $\beta=3$ ότι $\dfrac{-\alpha}{\beta}=\dfrac{\alpha}{-\beta}=-\dfrac{\alpha}{\beta}$.

\zoni{Για το σπίτι}{μοίρασε και σκέψου}

\ex Φτιάξε ένα πρόβλημα όπου ένα χρέος μοιράζεται ισόποσα και λύσε το.

\paragrafos{Δεκαδική μορφή ρητών: όταν τα ψηφία δεν τελειώνουν}{Α.7.7, «Δεκαδική μορφή ρητών αριθμών»}

\begin{anagke}
Ο Στέλιος παίζει με το κομπιουτεράκι του κινητού και πατάει $1:3$. Βλέπει $0{,}33333333$ και συνεχίζει με το $1:7$: $0{,}14285714$. «Το $1:7$ έχει ένα σωρό ψηφία!». Η Πολυξένη το αντιγράφει και προσπαθεί να δει αν υπάρχει κάποιο μοτίβο: «$142857\ 14\ldots$ Το $142857$ επαναλαμβάνεται!». Δοκιμάζει με το $2:7$ και βρίσκει $0{,}285714\ 285\ldots$: ίδια ψηφία, αλλά ξεκινούν από άλλο σημείο.

Ο Νίκος δοκιμάζει $3:8$ και βρίσκει $0{,}375$ και τελειώνει. Η Αγγελική ρωτάει: «Πότε τελειώνει η διαίρεση και πότε όχι;». Η Μυρτώ κάνει την ευκλείδεια διαίρεση στο χέρι: «Στο $1:7$ τα υπόλοιπα είναι πάντα $1,2,\dots,6$, δεν έχουν άλλες επιλογές. Μετά από έξι βήματα θα ξαναβρούμε κάποιο υπόλοιπο και τα ψηφία θα επαναληφθούν».

Και ο Βασίλης προτείνει το δύσκολο: «Αν πάρω το $0{,}3333\ldots$ και το πολλαπλασιάσω επί τρία, βγαίνει $0{,}9999\ldots$. Αλλά το $\frac{1}{3}\cdot3=1$. Άρα το $0{,}9999\ldots$ ισούται με το $1$;». Σε αυτή την παράγραφο η παρέα βρίσκει τον κανόνα του πότε τελειώνει ένα δεκαδικό ανάπτυγμα, μαθαίνει να γράφει τους άπειρους δεκαδικούς και πώς να γυρίζει πίσω στο κλάσμα.
\end{anagke}

\begin{lexi}{περιοδικός, περίοδος}{περιοδικός, περίοδος}
\lexiword{Περίοδος} από το «περί-οδος», δηλαδή «γύρω-δρόμος»: ένας δρόμος που τον ξανακάνουμε. Ο \lexiword{περιοδικός δεκαδικός} έχει ψηφία που επαναλαμβάνονται σε κύκλο (την περίοδο) για πάντα. Ο \lexiword{ακριβής δεκαδικός} τελειώνει, χωρίς να επαναλαμβάνεται τίποτα.
\end{lexi}

\begin{metafrasi}
Το $\frac{3}{8}$ τελειώνει. & $0{,}375$\\
Το $\frac{1}{3}$ επαναλαμβάνει το $3$. & $0{,}\overline{3}=0{,}333\ldots$\\
Το $\frac{1}{7}$ επαναλαμβάνει το $142857$. & $0{,}\overline{142857}$\\
Το $\frac{5}{6}$: το $8$ μία φορά και μετά το $3$. & $0{,}8\overline{3}=0{,}8333\ldots$\\
Από περιοδικό σε κλάσμα. & $0{,}\overline{27}=\dfrac{27}{99}=\dfrac{3}{11}$\\
\end{metafrasi}

\tmima{Η θεωρία}

Κάθε ρητός $\frac{\kappa}{\lambda}$ γίνεται δεκαδικός αν κάνουμε τη διαίρεση $\kappa:\lambda$. Υπάρχουν δύο ενδεχόμενα:

\begin{itemize}[leftmargin=1.2em,itemsep=0.3ex,topsep=2pt]
\item Η διαίρεση \textbf{τελειώνει} (υπόλοιπο $0$): έχουμε \textbf{ακριβή δεκαδικό} (πεπερασμένο πλήθος δεκαδικών ψηφίων). Π.χ. $\frac{3}{8}=0{,}375$, $\frac{9}{20}=0{,}45$.
\item Η διαίρεση \textbf{δεν τελειώνει}: τα υπόλοιπα είναι πάντα μικρότερα από τον διαιρέτη $\lambda$, άρα κάποτε επαναλαμβάνονται και μαζί τους επαναλαμβάνονται τα ψηφία του πηλίκου. Έχουμε \textbf{περιοδικό δεκαδικό}. Γράφουμε $\frac{5}{6}=0{,}8\overline{3}$ και $\frac{4}{11}=0{,}\overline{36}$.
\end{itemize}

Το τμήμα που επαναλαμβάνεται λέγεται \textbf{περίοδος}. Αν η περίοδος αρχίζει αμέσως μετά την υποδιαστολή, ο δεκαδικός λέγεται \emph{απλός περιοδικός} ($0{,}\overline{36}$), αλλιώς \emph{μικτός} ($0{,}8\overline{3}$).

\begin{kanonas}{πότε τελειώνει}
Ένα ανάγωγο κλάσμα γράφεται ως ακριβής δεκαδικός όταν ο παρονομαστής του έχει \textbf{μόνο παράγοντες $2$ και $5$} (γιατί τότε γίνεται δεκαδικό κλάσμα με παρονομαστή δύναμη του $10$). Αλλιώς είναι περιοδικός.
\end{kanonas}

Έτσι το $\frac{7}{20}$ ($20=2^2\cdot5$) τελειώνει: $0{,}35$. Το $\frac{5}{12}$ ($12=2^2\cdot3$) δεν τελειώνει: $0{,}41\overline{6}$.

\tmimasmall{Από περιοδικό σε κλάσμα}

Αν $x=0{,}\overline{3}$, τότε $10x=3{,}\overline{3}$ και αφαιρώντας $10x-x=3$, άρα $9x=3$ και $x=\frac{1}{3}$. Με την ίδια ιδέα:
\begin{itemize}[leftmargin=1.2em,itemsep=0.3ex,topsep=2pt]
\item \textbf{Απλός περιοδικός:} αριθμητής η περίοδος και παρονομαστής τόσα εννιάρια όσα τα ψηφία της περιόδου: $0{,}\overline{27}=\frac{27}{99}=\frac{3}{11}$, $0{,}\overline{7}=\frac{7}{9}$, $1{,}\overline{3}=1+\frac{3}{9}=\frac{4}{3}$.
\item \textbf{Μικτός:} $0{,}2\overline{3}$: με $x=0{,}2\overline{3}$, $10x=2{,}\overline{3}=2+\frac{1}{3}=\frac{7}{3}$, άρα $x=\frac{7}{30}$.
\end{itemize}
Κάθε περιοδικός δεκαδικός είναι ρητός, και κάθε ρητός είναι ακριβής ή περιοδικός δεκαδικός. Ο $0{,}\overline{9}=1$: με $x=0{,}\overline{9}$, $10x=9{,}\overline{9}$, $9x=9$, $x=1$.

\begin{prosoxi}
\begin{itemize}[leftmargin=1.2em,itemsep=0.3ex,topsep=2pt]
\item Το $\frac{1}{3}$ δεν είναι ίσο με $0{,}33$ αλλά με $0{,}\overline{3}$: το $0{,}33$ είναι προσέγγιση.
\item Η περίοδος είναι μόνο το κομμάτι που επαναλαμβάνεται: στο $0{,}8\overline{3}$ το $8$ δεν επαναλαμβάνεται.
\item Η περίοδος ενός κλάσματος με παρονομαστή $\lambda$ έχει το πολύ $\lambda-1$ ψηφία (λόγω των υπολοίπων).
\item Τα κομπιουτεράκια κόβουν ή στρογγυλοποιούν τα ψηφία: $1:7$ δεν σταματά στο $0{,}14285714$.
\end{itemize}
\end{prosoxi}

\tmima{Λυμένα παραδείγματα}

\paradeigma{Ακριβής ή περιοδικός;}
\tixero Τα κλάσματα $\frac{7}{20}$, $\frac{5}{12}$, $\frac{9}{14}$, $\frac{11}{40}$.\par
\tithelo Ποια έχουν ακριβή δεκαδικό.\par
\possindeo Αναλύω τους παρονομαστές: $20=2^2\cdot5$, $12=2^2\cdot3$, $14=2\cdot7$, $40=2^3\cdot5$. Μόνο παράγοντες $2$ και $5$ έχουν το $20$ και το $40$.\par
\apantisi Ακριβείς: $\frac{7}{20}=0{,}35$ και $\frac{11}{40}=0{,}275$. Περιοδικοί: $\frac{5}{12}=0{,}41\overline{6}$ και $\frac{9}{14}=0{,}6\overline{428571}$.\par
\elegxos Το $0{,}35=\frac{35}{100}=\frac{7}{20}$.

\paradeigma{Η διαίρεση του 1:7}
\tixero Το $\frac{1}{7}$.\par
\tithelo Να βρω τα ψηφία και την περίοδο.\par
\possindeo $1:7$: $10:7=1$ υπόλοιπο $3$, $30:7=4$ υπόλοιπο $2$, $20:7=2$ υπόλοιπο $6$, $60:7=8$ υπόλοιπο $4$, $40:7=5$ υπόλοιπο $5$, $50:7=7$ υπόλοιπο $1$. Το υπόλοιπο $1$ ξαναεμφανίζεται, άρα η περίοδος είναι $142857$.\par
\apantisi $\frac{1}{7}=0{,}\overline{142857}$.\par
\elegxos Η περίοδος έχει $6$ ψηφία, όσα τα δυνατά υπόλοιπα $1,\dots,6$.

\paradeigma{Από περιοδικό σε κλάσμα}
\tixero Οι $0{,}\overline{27}$, $2{,}\overline{6}$ και $0{,}1\overline{6}$.\par
\tithelo Να τους γράψω ως κλάσματα.\par
\possindeo $0{,}\overline{27}=\frac{27}{99}=\frac{3}{11}$. $2{,}\overline{6}=2+\frac{6}{9}=2+\frac{2}{3}=\frac{8}{3}$. Για τον τρίτο $x=0{,}1\overline{6}$, $10x=1{,}\overline{6}=1+\frac{2}{3}=\frac{5}{3}$, άρα $x=\frac{5}{30}=\frac{1}{6}$.\par
\apantisi $\frac{3}{11}$, $\frac{8}{3}$, $\frac{1}{6}$.\par
\elegxos Το $\frac{1}{6}=0{,}1666\ldots$ και το $3:11=0{,}2727\ldots$.

\paradeigma{Στρογγυλοποίηση περιοδικών}
\tixero Οι $\frac{2}{3}$, $\frac{5}{6}$, $\frac{1}{7}$.\par
\tithelo Στρογγυλοποίησε στα εκατοστά.\par
\possindeo $\frac{2}{3}=0{,}666\ldots\approx0{,}67$. $\frac{5}{6}=0{,}8333\ldots\approx0{,}83$. $\frac{1}{7}=0{,}1428\ldots\approx0{,}14$.\par
\apantisi $0{,}67$, $0{,}83$, $0{,}14$.\par
\elegxos Το σφάλμα στο πρώτο: $0{,}67-0{,}666\ldots=0{,}0033\ldots<0{,}005$.

\paradeigma{Ισούται το $0{,}\overline{9}$ με το $1$;}
\tixero Ο αριθμός $0{,}\overline{9}$.\par
\tithelo Να τον συγκρίνω με το $1$.\par
\possindeo Αν $x=0{,}\overline{9}$, τότε $10x=9{,}\overline{9}$. Αφαιρώ: $10x-x=9$, άρα $9x=9$ και $x=1$. Ή: $\frac{1}{3}=0{,}\overline{3}$, άρα $3\cdot\frac{1}{3}=3\cdot0{,}\overline{3}=0{,}\overline{9}$, δηλαδή $1=0{,}\overline{9}$.\par
\apantisi $0{,}\overline{9}=1$.\par
\elegxos Δεν υπάρχει αριθμός ανάμεσα στο $0{,}\overline{9}$ και στο $1$.

\tmima{Ασκήσεις}

\zoni{Το κατάλαβα}{ακριβείς και περιοδικοί}

\ex Γράψε σε δεκαδική μορφή (ακριβή ή περιοδική): α) $\frac{3}{8}$ \quad β) $\frac{5}{6}$ \quad γ) $\frac{4}{11}$ \quad δ) $\frac{7}{9}$ \quad ε) $\frac{9}{20}$ \quad στ) $\frac{2}{15}$ \quad ζ) $\frac{1}{12}$.

\ex Ποια από τα $\frac{7}{20}$, $\frac{5}{12}$, $\frac{9}{14}$, $\frac{11}{40}$, $\frac{3}{16}$, $\frac{13}{25}$ γράφονται ως ακριβείς δεκαδικοί; Χωρίς να κάνεις τη διαίρεση.

\ex Γράψε τα πρώτα οκτώ δεκαδικά ψηφία των $0{,}\overline{36}$, $2{,}1\overline{45}$ και $0{,}\overline{142857}$.

\ex Γράψε ως κλάσματα: α) $0{,}\overline{7}$ \quad β) $0{,}\overline{15}$ \quad γ) $1{,}\overline{3}$ \quad δ) $0{,}\overline{81}$ \quad ε) $0{,}2\overline{3}$.

\ex Σύγκρινε: α) $0{,}\overline{3}$ και $0{,}33$ \quad β) $0{,}\overline{6}$ και $\frac{2}{3}$ \quad γ) $0{,}7\overline{2}$ και $0{,}72$.

\ex Στρογγυλοποίησε στα εκατοστά: α) $\frac{2}{3}$ \quad β) $\frac{5}{6}$ \quad γ) $\frac{1}{7}$ \quad δ) $\frac{4}{9}$.

\ex Σωστό ή λάθος; α) Κάθε ρητός είναι ακριβής ή περιοδικός δεκαδικός. \quad β) $\frac{1}{3}=0{,}33$. \quad γ) $0{,}\overline{9}=1$. \quad δ) Το $\frac{1}{7}$ έχει περίοδο $142857$. \quad ε) Τα κλάσματα με παρονομαστή $8$ είναι ακριβείς δεκαδικοί. \quad στ) $\frac{1}{6}=0{,}16$.

\zoni{Το χρησιμοποιώ}{κομπιουτεράκι, μοιρασιές, μέσοι όροι}

\ex \textbf{Το κομπιουτεράκι.} Ο Στέλιος γράφει $1:7$ και βλέπει $0{,}14285714$. Ποια είναι η περίοδος; Βρες με το κομπιουτεράκι το $2:7$ και το $3:7$ και πες τι παρατηρείς.

\ex \textbf{Η μοιρασιά.} Τα $10$ € μοιράζονται σε $3$ παιδιά. Τι δείχνει το $10:3$ και πόσο παίρνει στην πράξη το καθένα;

\ex \textbf{Η ταχύτητα.} Ο Νίκος τρέχει $100$ μ. σε $12$ s. Πόση είναι η ταχύτητά του σε m/s ως περιοδικός και ως στρογγυλοποιημένος στα εκατοστά;

\ex \textbf{Τα ποσοστά.} Γράψε το $\frac{1}{3}$ και το $\frac{5}{6}$ ως ποσοστά με περιοδικούς και στρογγυλοποίησε στο δέκατο.

\ex \textbf{Ο μέσος όρος.} Οι βαθμοί $7$, $8$, $8$ έχουν μέσο όρο $\frac{23}{3}$. Γράψε τον ως περιοδικό και στρογγυλοποίησέ τον στα εκατοστά.

\zoni{Το σκέφτομαι}{γιατί επαναλαμβάνονται}

\ex Εξήγησε γιατί η διαίρεση $\kappa:7$ έχει περίοδο το πολύ $6$ ψηφίων.

\ex Δείξε ότι $0{,}\overline{9}=1$ με τον τρόπο $x=0{,}\overline{9}$, $10x=9{,}\overline{9}$.

\ex Βρες τη δεκαδική μορφή του $\frac{1}{11}$ και του $\frac{1}{13}$ (δοκίμασε με κομπιουτεράκι και με διαίρεση).

\ex Βρες έναν ρητό ανάμεσα στους $0{,}\overline{3}$ και $0{,}34$.

\ex Γιατί το $\frac{1}{3}$ ισούται με $0{,}\overline{3}$ και όχι με $0{,}333$; Τι μας λέει το $3\cdot0{,}\overline{3}$;

\zoni{Για το σπίτι}{ανακάλυψε τα μοτίβα}

\ex Υπολόγισε με κομπιουτεράκι το $1:\nu$ για $\nu=2,\dots,12$ και βρες ποια έχουν ακριβές ανάπτυγμα. Τι κοινό έχουν οι παρονομαστές τους;

\paragrafos{Δυνάμεις με εκθέτη φυσικό: όταν οι πολλαπλασιασμοί συνεχίζονται}{Α.7.8, «Δυνάμεις ρητών αριθμών με εκθέτη φυσικό»}

\begin{anagke}
Στο διάλειμμα η Πολυξένη λέει στην Αγγελική ένα μυστικό και της ζητάει να το πει μέσα σε μια ώρα σε δύο άλλα παιδιά. Κάθε παιδί που το μαθαίνει το λέει, μέσα στην επόμενη ώρα, σε δύο άλλα. «Μέσα σε δέκα ώρες θα το ξέρει όλο το σχολείο;» ρωτάει ο Ηλίας. Ο Στέλιος γράφει: $1,\ 2,\ 4,\ 8,\ 16,\ \ldots$ και κάθε φορά το πλήθος διπλασιάζεται. «Μετά από $10$ ώρες θα είναι $2\cdot2\cdot2\cdot\ldots$ δέκα φορές».

Ο Βασίλης λέει ότι $2$ δέκα φορές επί $2$ είναι $2\cdot10=20$. Ο Νίκος τον διορθώνει: «Όχι, δέκα φορές πολλαπλασιασμός, όχι δέκα φορές πρόσθεση. Θα βγει $1\,024$!». Η Ρήνη προτείνει μια σύντομη γραφή: $2\cdot2\cdot2\cdot2\cdot2=2^5$, το «δύο στην πέμπτη», που το ξέρουν από την Ε΄ και ΣΤ΄ Δημοτικού για τους φυσικούς. Η κυρία Αλεξάνδρα το βάζει στο νέο επίπεδο: «Μπορούμε να υψώσουμε στη δύναμη και έναν ρητό. Τι νομίζετε ότι είναι το $(-2)^3$; Και το $\left(-\frac{1}{2}\right)^3$;».

Σε αυτή την παράγραφο η παρέα ξαναβρίσκει τις δυνάμεις, μαθαίνει το πρόσημό τους και τις ιδιότητες που τις κάνουν γρήγορες, και προσέχει τα σημεία που μπερδεύουν τους περισσότερους: το $(-3)^2$ και το $-3^2$.
\end{anagke}

\begin{lexi}{δύναμη, βάση, εκθέτης}{δύναμη, βάση, εκθέτης}
Η \lexiword{δύναμη} είναι «ύψωση», δηλαδή το γινόμενο ίσων παραγόντων· στα αρχαία ελληνικά «δύναμις» ήταν και το τετράγωνο. Η \lexiword{βάση} είναι ο αριθμός που πολλαπλασιάζεται με τον εαυτό του και ο \lexiword{εκθέτης} (από το «εκ-τίθημι», «βάζω έξω») αυτός που γράφουμε πάνω δεξιά, δείχνοντας πόσους παράγοντες έχουμε.
\end{lexi}

\begin{metafrasi}
Πέντε φορές ο $2$ πολλαπλασιασμένος. & $2^5=2\cdot2\cdot2\cdot2\cdot2=32$\\
Ο $-2$ τρεις φορές. & $(-2)^3=(-2)(-2)(-2)=-8$\\
Μετά από ν γύρους το μυστικό το ξέρουν. & $2^\nu$ παιδιά\\
Πολλαπλασιάζω δυνάμεις με ίδια βάση. & $\alpha^\mu\cdot\alpha^\nu=\alpha^{\mu+\nu}$\\
Η δύναμη μιας δύναμης. & $(\alpha^\mu)^\nu=\alpha^{\mu\nu}$\\
\end{metafrasi}

\tmima{Η θεωρία}

\begin{orismos}{δύναμη με εκθέτη φυσικό}
Για κάθε ρητό $\alpha$ και φυσικό $\nu\ge2$ η \textbf{δύναμη} $\alpha^\nu$ είναι το γινόμενο $\nu$ παραγόντων ίσων με $\alpha$: $\alpha^\nu=\alpha\cdot\alpha\cdots\alpha$. Ορίζουμε ακόμη $\alpha^1=\alpha$.
\end{orismos}

Για παράδειγμα $(0{,}1)^3=0{,}1\cdot0{,}1\cdot0{,}1=0{,}001$ και $\left(\frac{2}{3}\right)^3=\frac{2}{3}\cdot\frac{2}{3}\cdot\frac{2}{3}=\frac{8}{27}$.

\tmimasmall{Το πρόσημο της δύναμης}

Από τον κανόνα των προσήμων του πολλαπλασιασμού:
\begin{itemize}[leftmargin=1.2em,itemsep=0.3ex,topsep=2pt]
\item Μια δύναμη με θετική βάση είναι πάντα θετική.
\item Μια δύναμη με αρνητική βάση είναι \textbf{θετική} αν ο εκθέτης είναι \textbf{άρτιος} και \textbf{αρνητική} αν είναι \textbf{περιττός}: $(-2)^4=16$, $(-2)^3=-8$, $(-1)^{100}=1$, $(-1)^{101}=-1$.
\end{itemize}

Η βάση είναι ό,τι βρίσκεται μέσα στην παρένθεση: $(-3)^2=9$ ενώ $-3^2=-(3\cdot3)=-9$.

\begin{idiotita}{ιδιότητες των δυνάμεων (εκθέτης φυσικός)}
Για ρητούς $\alpha,\,\beta$ και φυσικούς $\mu,\,\nu$ ισχύουν:
\begin{itemize}[leftmargin=1.2em,itemsep=0.3ex,topsep=2pt]
\item $\alpha^\mu\cdot\alpha^\nu=\alpha^{\mu+\nu}$ \qquad $\alpha^\mu:\alpha^\nu=\alpha^{\mu-\nu}$ (με $\alpha\neq0$, $\mu>\nu$)
\item $(\alpha\cdot\beta)^\nu=\alpha^\nu\cdot\beta^\nu$ \qquad $\left(\dfrac{\alpha}{\beta}\right)^\nu=\dfrac{\alpha^\nu}{\beta^\nu}$ (με $\beta\neq0$)
\item $(\alpha^\mu)^\nu=\alpha^{\mu\cdot\nu}$
\end{itemize}
\end{idiotita}

Γιατί; Το $2^3\cdot2^4=(2\cdot2\cdot2)\cdot(2\cdot2\cdot2\cdot2)=2^7$ (τρεις παράγοντες και τέσσερις: επτά). Το $(3^2)^3=3^2\cdot3^2\cdot3^2=3^6$. Το $(2\cdot5)^3=2\cdot5\cdot2\cdot5\cdot2\cdot5=2^3\cdot5^3=8\cdot125=1\,000$.

\tmimasmall{Προτεραιότητα πράξεων}

Η σειρά είναι: παρενθέσεις, μετά δυνάμεις, μετά πολλαπλασιασμοί και διαιρέσεις, τέλος προσθέσεις και αφαιρέσεις. Για παράδειγμα $2\cdot3^2-(-2)^3+10:(-5)=18-(-8)+(-2)=18+8-2=24$.

\begin{prosoxi}
\begin{itemize}[leftmargin=1.2em,itemsep=0.3ex,topsep=2pt]
\item $2^3=8$, όχι $6$: η δύναμη είναι επανάληψη πολλαπλασιασμού, όχι γινόμενο $2\cdot3$.
\item $(-3)^2=+9$, αλλά $-3^2=-9$ (ο εκθέτης ενεργεί μόνο στο $3$).
\item Το $(\alpha+\beta)^2$ δεν είναι $\alpha^2+\beta^2$: $(3+4)^2=49$ ενώ $3^2+4^2=25$.
\item Οι ιδιότητες $\alpha^\mu\cdot\alpha^\nu=\alpha^{\mu+\nu}$ ισχύουν μόνο για \emph{ίδια βάση}: το $2^3\cdot3^2$ δεν γίνεται $6^5$.
\item Στο $\alpha^\mu\cdot\alpha^\nu$ προσθέτω εκθέτες, δεν πολλαπλασιάζω: $2^3\cdot2^4=2^7$ και όχι $2^{12}$.
\end{itemize}
\end{prosoxi}

\tmima{Λυμένα παραδείγματα}

\paradeigma{Δυνάμεις και πρόσημα}
\tixero Οι δυνάμεις $2^5$, $(-3)^2$, $-3^2$, $(-3)^3$, $\left(-\frac{1}{2}\right)^3$, $(0{,}1)^3$.\par
\tithelo Να τις υπολογίσω.\par
\possindeo $2^5=32$. $(-3)^2=9$. $-3^2=-(3\cdot3)=-9$. $(-3)^3=-27$. $\left(-\frac{1}{2}\right)^3=-\frac{1}{8}$. $(0{,}1)^3=0{,}001$.\par
\apantisi $32$, $9$, $-9$, $-27$, $-\frac{1}{8}$, $0{,}001$.\par
\elegxos Αρνητική βάση με περιττό εκθέτη δίνει αρνητικό: το $(-3)^3$ και το $\left(-\frac{1}{2}\right)^3$.

\paradeigma{Με τις ιδιότητες}
\tixero Οι παραστάσεις $2^3\cdot2^4$, $5^8:5^5$, $(3^2)^3$, $(2\cdot5)^3$, $\left(\frac{2}{3}\right)^3$.\par
\tithelo Να τις υπολογίσω.\par
\possindeo $2^{3+4}=2^7=128$. $5^{8-5}=5^3=125$. $3^{2\cdot3}=3^6=729$. $2^3\cdot5^3=8\cdot125=1\,000$. $\frac{2^3}{3^3}=\frac{8}{27}$.\par
\apantisi $128$, $125$, $729$, $1\,000$, $\frac{8}{27}$.\par
\elegxos $(2\cdot5)^3=10^3=1\,000$.

\paradeigma{Το μυστικό του σχολείου}
\tixero Κάθε παιδί λέει το μυστικό σε $2$ νέα παιδιά κάθε ώρα. Στην αρχή το ξέρει $1$ παιδί.\par
\tithelo Πόσα παιδιά το ξέρουν μετά από $10$ ώρες.\par
\possindeo Μετά από $1$ ώρα: $2$. Μετά από $2$: $4$. Μετά από $3$: $8$. Μετά από $\nu$ ώρες το πλήθος διπλασιάζεται $\nu$ φορές: $2^\nu$. Για $\nu=10$: $2^{10}=1\,024$.\par
\apantisi $1\,024$ παιδιά.\par
\elegxos Από την $9η$ ώρα ($512$) στη $10η$ διπλασιάζεται: $1\,024$.

\paradeigma{Η προτεραιότητα}
\tixero Η παράσταση $2\cdot3^2-(-2)^3+10:(-5)$.\par
\tithelo Να την υπολογίσω.\par
\possindeo Πρώτα οι δυνάμεις: $3^2=9$, $(-2)^3=-8$. Μετά πολλαπλασιασμοί και διαιρέσεις: $2\cdot9=18$, $10:(-5)=-2$. Τέλος $18-(-8)+(-2)=18+8-2=24$.\par
\apantisi $24$.\par
\elegxos Αν έκανα πρώτα τον πολλαπλασιασμό $2\cdot3=6$, θα έβγαινε $6^2=36$: λάθος σειρά.

\tmima{Ασκήσεις}

\zoni{Το κατάλαβα}{δυνάμεις, πρόσημα, ιδιότητες}

\ex Υπολόγισε: α) $2^5$ \quad β) $3^4$ \quad γ) $(-2)^3$ \quad δ) $(-2)^4$ \quad ε) $(-1)^7$ \quad στ) $0^5$.

\ex Υπολόγισε: α) $(-3)^2$ \quad β) $-3^2$ \quad γ) $\left(-\frac{1}{2}\right)^3$ \quad δ) $(0{,}1)^3$ \quad ε) $(-0{,}2)^2$ \quad στ) $\left(\frac{2}{3}\right)^3$.

\ex Πες το πρόσημο χωρίς να υπολογίσεις: α) $(-5)^6$ \quad β) $(-5)^7$ \quad γ) $-(-3)^4$ \quad δ) $(-2)^3\cdot(-3)^2$.

\ex Γράψε ως μία δύναμη: α) $2^3\cdot2^4$ \quad β) $5^8:5^5$ \quad γ) $(3^2)^3$ \quad δ) $7^2\cdot7\cdot7^3$ \quad ε) $(-2)^5\cdot(-2)^2$ \quad στ) $(-4)^6:(-4)^4$.

\ex Υπολόγισε με τις ιδιότητες: α) $(2\cdot5)^3$ \quad β) $(-2\cdot3)^2$ \quad γ) $\left(\frac{2}{3}\right)^3$ \quad δ) $\left(-\frac{3}{4}\right)^2$.

\ex Υπολόγισε: α) $2\cdot3^2-(-2)^3+10:(-5)$ \quad β) $(-2)^2\cdot3-4^2:(-2)$ \quad γ) $5-3\cdot(-2)^2$ \quad δ) $\big((-1)^5+(-1)^6\big)\cdot10$.

\ex Σωστό ή λάθος; α) $2^3=6$. \quad β) $(-2)^4=16$. \quad γ) $-2^4=16$. \quad δ) $(\alpha\beta)^2=\alpha^2\beta^2$. \quad ε) $\alpha^2\cdot\alpha^3=\alpha^6$. \quad στ) $(\alpha^2)^3=\alpha^6$. \quad ζ) $(\alpha+\beta)^2=\alpha^2+\beta^2$.

\zoni{Το χρησιμοποιώ}{φήμες, κύβοι, σκάκι}

\ex \textbf{Το μυστικό.} Κάθε παιδί λέει το μυστικό σε $2$ νέα κάθε ώρα. Φτιάξε πίνακα με το πλήθος των παιδιών μετά από $0,\,1,\,\dots,\,10$ ώρες και πες πότε ξεπερνά τα $1\,000$.

\ex \textbf{Ο κύβος.} Ένας κύβος έχει ακμή $12$ cm. Βρες τον όγκο ($12^3$) και το εμβαδόν της επιφάνειάς του ($6\cdot12^2$).

\ex \textbf{Το σκάκι.} Στο πρώτο τετράγωνο ενός σκακιστικού πίνακα βάζουμε $1$ κόκκο, στο δεύτερο $2$, στο τρίτο $4$ και συνεχίζουμε διπλασιάζοντας. Πόσοι κόκκοι είναι στο $8$ο τετράγωνο και πόσοι στο $11$ο;

\ex \textbf{Τα βακτήρια.} Ένα βακτήριο τριπλασιάζεται κάθε ώρα. Πόσα είναι μετά από $5$ ώρες;

\ex \textbf{Οι τετράγωνοι αριθμοί.} Οι αριθμοί $1,\,4,\,9,\,16,\dots$ είναι τα τετράγωνα. Ποιος είναι ο $20$ός και ο $\nu$-οστός;

\ex \textbf{Η τιμή διπλασιάζεται.} Μια τιμή των $3$ € διπλασιάζεται $4$ φορές. Πόσο γίνεται;

\ex \textbf{Η δύναμη του $10$.} Υπολόγισε $10^3\cdot10^4$ και $10^8:10^5$.

\zoni{Το σκέφτομαι}{όταν το πρόσημο και ο εκθέτης παίζουν}

\ex Ο Βασίλης λέει ότι $2^3=6$. Βρες το λάθος και εξήγησε τη διαφορά ανάμεσα στο $2\cdot3$ και στο $2^3$.

\ex Δείξε με $\alpha=3$ και $\beta=4$ ότι $(\alpha+\beta)^2\neq\alpha^2+\beta^2$.

\ex Πότε ισχύει $(-\alpha)^\nu=\alpha^\nu$ και πότε $(-\alpha)^\nu=-\alpha^\nu$;

\ex Πόσο είναι τα $1^\nu$, $0^\nu$ (για $\nu\ge1$) και $(-1)^\nu$;

\ex Εξήγησε γιατί $2^3\cdot2^4=2^7$ γράφοντας όλους τους παράγοντες.

\ex Ποιο είναι μεγαλύτερο: το $3^4$ ή το $4^3$;

\zoni{Για το σπίτι}{δίπλωσε και μέτρα}

\ex Δίπλωσε ένα χαρτί στη μέση όσες φορές μπορείς. Πόσα στρώματα έχεις μετά από $1,\,2,\,3,\,4,\,5$ διπλώματα;

\paragrafos{Δυνάμεις με εκθέτη ακέραιο: προς τα αριστερά του μηδενός}{Α.7.9, «Δυνάμεις ρητών αριθμών με εκθέτη ακέραιο»}

\begin{anagke}
Η Πολυξένη γράφει στον πίνακα τις δυνάμεις του $2$: $2^4=16$, $2^3=8$, $2^2=4$, $2^1=2$. Ο Στέλιος κοιτάζει και λέει: «Κάθε φορά που ο εκθέτης μειώνεται κατά $1$, ο αριθμός γίνεται το μισό: $16$, $8$, $4$, $2$». «Και αν συνεχίσουμε;» ρωτάει η Αγγελική. Ο επόμενος αριθμός πρέπει να είναι το μισό του $2$, δηλαδή $1$: $2^0=1$. Και μετά το μισό του $1$: $\frac{1}{2}$, δηλαδή $2^{-1}=\frac{1}{2}$. Και $2^{-2}=\frac{1}{4}$, $2^{-3}=\frac{1}{8}$.

«Αρνητικός εκθέτης;» αναρωτιέται ο Βασίλης. «Δεν μπορώ να πολλαπλασιάσω τον $2$ μείον τρεις φορές!». Η κυρία Αλεξάνδρα εξηγεί: «Δεν χρειάζεται να το κάνετε. Αφού ο εκθέτης που μικραίνει διαιρεί με το $2$, μπορούμε να ορίσουμε ότι συνεχίζεται το ίδιο μοτίβο και στους αρνητικούς, ώστε να ισχύουν οι ιδιότητες των δυνάμεων που ξέρετε». Ο Νίκος το συνδέει με τις μονάδες: «Ένα χιλιοστό είναι το $\frac{1}{1\,000}$ του μέτρου, δηλαδή $10^{-3}$ μέτρα. Αυτό θα γράφουμε στο εργαστήριο».

Στην παράγραφο αυτή η παρέα αποκτά ένα εργαλείο για τους πολύ μεγάλους και τους πολύ μικρούς αριθμούς, και δίνει νόημα στο $\alpha^0$ και στο $\alpha^{-\nu}$.
\end{anagke}

\begin{lexi}{αρνητικός εκθέτης}{αντίστροφος, δύναμη}
Στα μαθηματικά οι νέες έννοιες ορίζονται ώστε να «συνεχίζουν το μοτίβο» και να ισχύουν οι παλιοί κανόνες. Η δύναμη με αρνητικό εκθέτη δεν σημαίνει «πολλαπλασιάζω αρνητικές φορές» αλλά «αντίστροφος της δύναμης»: το $2^{-3}$ είναι ο \lexiword{αντίστροφος} του $2^3$.
\end{lexi}

\begin{metafrasi}
Το μισό του $2$. & $2^0=1$\\
Το μισό του $1$. & $2^{-1}=\dfrac{1}{2}$\\
Ένα χιλιοστό του μέτρου. & $10^{-3}$ μ.\\
Η δύναμη με αρνητικό εκθέτη. & $\alpha^{-\nu}=\dfrac{1}{\alpha^{\nu}}$\\
Κλάσμα με αρνητικό εκθέτη. & $\left(\dfrac{\alpha}{\beta}\right)^{-\nu}=\left(\dfrac{\beta}{\alpha}\right)^{\nu}$\\
\end{metafrasi}

\tmima{Η θεωρία}

Ας γράψουμε τις δυνάμεις του $2$ και ας συνεχίσουμε προς τα δεξιά:
\begin{center}
\renewcommand{\arraystretch}{1.3}
\begin{tabular}{c|ccccccc}
$\nu$ & $3$ & $2$ & $1$ & $0$ & $-1$ & $-2$ & $-3$\\\hline
$2^\nu$ & $8$ & $4$ & $2$ & $1$ & $\frac{1}{2}$ & $\frac{1}{4}$ & $\frac{1}{8}$
\end{tabular}
\end{center}
Κάθε φορά που ο εκθέτης μειώνεται κατά $1$, η τιμή διαιρείται με $2$. Έτσι ορίζουμε:

\begin{orismos}{δύναμη με εκθέτη ακέραιο}
Για κάθε ρητό $\alpha\neq0$ και φυσικό $\nu$:
$$\alpha^0=1,\qquad \alpha^{-\nu}=\dfrac{1}{\alpha^{\nu}}.$$
\end{orismos}

Για παράδειγμα $5^0=1$, $(-7)^0=1$, $3^{-2}=\frac{1}{9}$, $(-2)^{-3}=\frac{1}{(-2)^3}=-\frac{1}{8}$, $10^{-3}=\frac{1}{1\,000}=0{,}001$. Το $0^0$ και τα $0^{-\nu}$ δεν ορίζονται.

\tmimasmall{Κλάσμα στη δύναμη}

Αν η βάση είναι κλάσμα, ο αρνητικός εκθέτης αναποδογυρίζει το κλάσμα:
$$\left(\dfrac{\alpha}{\beta}\right)^{-\nu}=\left(\dfrac{\beta}{\alpha}\right)^{\nu}\quad(\alpha,\beta\neq0).$$
Για παράδειγμα $\left(\frac{1}{2}\right)^{-3}=2^3=8$ και $\left(\frac{2}{3}\right)^{-2}=\left(\frac{3}{2}\right)^2=\frac{9}{4}$.

\begin{idiotita}{ιδιότητες με ακέραιο εκθέτη}
Οι ιδιότητες των δυνάμεων ισχύουν και για ακέραιους εκθέτες ($\alpha,\beta\neq0$): $\alpha^\mu\cdot\alpha^\nu=\alpha^{\mu+\nu}$, $\alpha^\mu:\alpha^\nu=\alpha^{\mu-\nu}$, $(\alpha\beta)^\nu=\alpha^\nu\beta^\nu$, $\left(\frac{\alpha}{\beta}\right)^\nu=\frac{\alpha^\nu}{\beta^\nu}$, $(\alpha^\mu)^\nu=\alpha^{\mu\nu}$.
\end{idiotita}

Με το νέο ορισμό η διαίρεση δεν χρειάζεται περιορισμό $\mu>\nu$: $5^3:5^5=5^{-2}=\frac{1}{25}$. Και βγαίνει τώρα φυσιολογικά ότι $\alpha^0=1$: $\alpha^3:\alpha^3=\alpha^{3-3}=\alpha^0$ και, αφού $\alpha^3:\alpha^3=1$, είναι $\alpha^0=1$.

\begin{prosoxi}
\begin{itemize}[leftmargin=1.2em,itemsep=0.3ex,topsep=2pt]
\item Το $2^{-3}$ δεν είναι $-8$ ούτε $-6$: είναι $\frac{1}{8}$, θετικό. Το αρνητικό σημείο στον εκθέτη δεν αλλάζει το πρόσημο.
\item $\alpha^0=1$ και όχι $0$ (για $\alpha\neq0$).
\item $\left(\frac{1}{2}\right)^{-3}=8$, όχι $\frac{1}{8}$: αναποδογυρίζω το κλάσμα και κάνω τον εκθέτη θετικό.
\item Το πρόσημο της βάσης παίζει ρόλο: $(-2)^{-3}=-\frac{1}{8}$ και $(-2)^{-2}=\frac{1}{4}$.
\item Το $0^{-1}$ δεν ορίζεται (θα σήμαινε $\frac{1}{0}$).
\end{itemize}
\end{prosoxi}

\tmima{Λυμένα παραδείγματα}

\paradeigma{Συνεχίζω τον πίνακα}
\tixero Οι δυνάμεις του $3$: $3^3=27$, $3^2=9$, $3^1=3$.\par
\tithelo Να βρω τα $3^0$, $3^{-1}$, $3^{-2}$, $3^{-3}$.\par
\possindeo Κάθε φορά ο εκθέτης μειώνεται και η τιμή διαιρείται με $3$: $3:3=1$, $1:3=\frac{1}{3}$, $\frac{1}{3}:3=\frac{1}{9}$, $\frac{1}{9}:3=\frac{1}{27}$.\par
\apantisi $3^0=1$, $3^{-1}=\frac{1}{3}$, $3^{-2}=\frac{1}{9}$, $3^{-3}=\frac{1}{27}$.\par
\elegxos Σύμφωνα με τον ορισμό $3^{-2}=\frac{1}{3^2}=\frac{1}{9}$.

\paradeigma{Υπολογισμοί}
\tixero Οι $3^{-2}$, $(-2)^{-3}$, $\left(\frac{1}{2}\right)^{-3}$, $\left(\frac{2}{3}\right)^{-2}$, $5^0$.\par
\tithelo Να τους υπολογίσω.\par
\possindeo $3^{-2}=\frac{1}{9}$. $(-2)^{-3}=\frac{1}{(-2)^3}=-\frac{1}{8}$. $\left(\frac{1}{2}\right)^{-3}=2^3=8$. $\left(\frac{2}{3}\right)^{-2}=\left(\frac{3}{2}\right)^2=\frac{9}{4}$. $5^0=1$.\par
\apantisi $\frac{1}{9}$, $-\frac{1}{8}$, $8$, $\frac{9}{4}$, $1$.\par
\elegxos $\left(\frac{1}{2}\right)^{-3}\cdot\left(\frac{1}{2}\right)^{3}=\left(\frac{1}{2}\right)^0=1$ και $8\cdot\frac{1}{8}=1$.

\paradeigma{Με τις ιδιότητες}
\tixero Οι $2^{-3}\cdot2^5$, $10^4:10^{-2}$, $(3^{-1})^2$, $(2\cdot5)^{-2}$.\par
\tithelo Να τους υπολογίσω.\par
\possindeo $2^{-3+5}=2^2=4$. $10^{4-(-2)}=10^6$. $3^{-2}=\frac{1}{9}$. $2^{-2}\cdot5^{-2}=\frac{1}{4}\cdot\frac{1}{25}=\frac{1}{100}$.\par
\apantisi $4$, $10^6=1\,000\,000$, $\frac{1}{9}$, $\frac{1}{100}$.\par
\elegxos $(2\cdot5)^{-2}=10^{-2}=\frac{1}{100}$.

\paradeigma{Παράσταση}
\tixero Η παράσταση $\left(\frac{1}{2}\right)^{-2}+3^0-2^{-1}$.\par
\tithelo Να την υπολογίσω.\par
\possindeo $\left(\frac{1}{2}\right)^{-2}=2^2=4$, $3^0=1$, $2^{-1}=\frac{1}{2}$. Άρα $4+1-\frac{1}{2}=\frac{9}{2}$.\par
\apantisi $\frac{9}{2}$.\par
\elegxos $4{,}5=4+1-0{,}5$.

\paradeigma{Εξίσωση με δύναμη}
\tixero Η εξίσωση $2^x=\frac{1}{16}$.\par
\tithelo Να βρω το $x$.\par
\possindeo Γράφω $\frac{1}{16}=\frac{1}{2^4}=2^{-4}$. Άρα $x=-4$.\par
\apantisi $x=-4$.\par
\elegxos $2^{-4}=\frac{1}{16}$.

\tmima{Ασκήσεις}

\zoni{Το κατάλαβα}{αρνητικοί εκθέτες και μηδέν}

\ex Υπολόγισε: α) $2^{-1}$ \quad β) $2^{-3}$ \quad γ) $10^{-2}$ \quad δ) $5^{-2}$ \quad ε) $4^0$ \quad στ) $(-3)^0$.

\ex Γράψε ως δεκαδικούς: α) $10^{-1}$ \quad β) $10^{-3}$ \quad γ) $10^{-5}$.

\ex Υπολόγισε: α) $(-2)^{-3}$ \quad β) $(-2)^{-2}$ \quad γ) $(-1)^{-5}$ \quad δ) $\left(\frac{1}{2}\right)^{-3}$ \quad ε) $\left(\frac{2}{3}\right)^{-2}$ \quad στ) $\left(-\frac{3}{4}\right)^{-1}$.

\ex Γράψε ως δύναμη με αρνητικό εκθέτη: α) $\frac{1}{9}$ (με βάση $3$) \quad β) $\frac{1}{32}$ (βάση $2$) \quad γ) $0{,}0001$ (βάση $10$) \quad δ) $\frac{1}{(-5)^3}$.

\ex Υπολόγισε με τις ιδιότητες: α) $2^{-3}\cdot2^5$ \quad β) $10^4:10^{-2}$ \quad γ) $(3^{-1})^2$ \quad δ) $(2\cdot5)^{-2}$ \quad ε) $7^{-4}\cdot7^4$ \quad στ) $(-2)^{-3}\cdot(-2)^5$.

\ex Υπολόγισε: α) $\left(\frac{1}{2}\right)^{-2}+3^0-2^{-1}$ \quad β) $2^{-2}+2^{-1}$ \quad γ) $10^2\cdot10^{-5}\cdot10^3$ \quad δ) $(3^2)^{-1}$ \quad ε) $5^0\cdot(-2)^{-2}$.

\ex Βρες το $x$: α) $2^x=\frac{1}{16}$ \quad β) $3^x=\frac{1}{27}$ \quad γ) $10^x=0{,}001$ \quad δ) $\left(\frac{1}{2}\right)^x=8$ \quad ε) $5^x=1$.

\ex Σωστό ή λάθος; α) $2^{-1}=-2$. \quad β) $\alpha^0=1$ για $\alpha\neq0$. \quad γ) $\left(\frac{1}{2}\right)^{-1}=2$. \quad δ) $10^{-3}=0{,}003$. \quad ε) $\alpha^{-\nu}=\frac{1}{\alpha^\nu}$. \quad στ) $2^{-2}\cdot2^2=1$. \quad ζ) $(-2)^{-2}=-\frac{1}{4}$.

\zoni{Το χρησιμοποιώ}{μονάδες, μισά, αξίες}

\ex \textbf{Οι μονάδες.} Γράψε με δυνάμεις του $10$ σε μέτρα: $1$ χιλιόμετρο, $1$ εκατοστό, $1$ χιλιοστό, $1$ μικρόμετρο ($0{,}000001$ μ.), $1$ νανόμετρο ($0{,}000000001$ μ.).

\ex \textbf{Η πίτα.} Κόβουμε κάθε φορά το μισό από ό,τι έμεινε. Τι μέρος της πίτας έχει μείνει μετά από $5$ κοψίματα; Γράψε το ως δύναμη του $2$.

\ex \textbf{Το ραδιενεργό υλικό.} Ένα υλικό 64 g χάνει το μισό της μάζας του κάθε ημέρα. Πόσο θα είναι μετά από $3$ ημέρες; Πόσο ήταν πριν από $2$ ημέρες;

\ex \textbf{Ο πίνακας.} Συμπλήρωσε τις δυνάμεις του $3$ από $3^4$ έως $3^{-3}$.

\ex \textbf{Το κινητό.} Η αξία ενός κινητού των $400$ € μειώνεται στο μισό κάθε χρόνο. Πόση αξία έχει μετά από $3$ έτη; Γράψε το με δύναμη.

\zoni{Το σκέφτομαι}{γιατί το μηδέν δίνει ένα}

\ex Γιατί $\alpha^0=1$ και όχι $0$; Χρησιμοποίησε το μοτίβο των δυνάμεων του $2$.

\ex Δείξε ότι $\left(\frac{2}{3}\right)^{-2}=\left(\frac{3}{2}\right)^2$.

\ex Σε τι διαφέρουν τα $2^{-3}$ και $-2^3$;

\ex Αν ισχύει $\alpha^\mu\cdot\alpha^\nu=\alpha^{\mu+\nu}$ και $\mu=-3$, $\nu=3$, τι σημαίνει $\alpha^0$;

\ex Γιατί δεν ορίζεται το $0^{-1}$;

\ex Υπολόγισε τα $(-1)^{-1}$, $(-1)^{-2}$, $(-1)^{-3}$ και βρες τον κανόνα.

\zoni{Για το σπίτι}{πόσο λεπτό είναι ένα φύλλο;}

\ex Ένα φύλλο χαρτιού έχει πάχος περίπου $0{,}1$ mm. Γράψε το πάχος σε μέτρα με δύναμη του $10$.

\paragrafos{Τυποποιημένη μορφή μεγάλων και μικρών αριθμών}{Α.7.10, «Τυποποιημένη μορφή μεγάλων και μικρών αριθμών»}

\begin{anagke}
Η Α3 επισκέπτεται το εργαστήριο βιολογίας του σχολείου και η κυρία Αλεξάνδρα δείχνει ένα μικροσκόπιο. Κάτω από τον φακό υπάρχουν μια τρίχα, ένα ερυθρό αιμοσφαίριο και ένα βακτήριο. Η καθηγήτρια γράφει στον πίνακα τα μεγέθη: η τρίχα έχει πάχος $0{,}000\,07$ μ., το αιμοσφαίριο $0{,}000\,007$ μ., το βακτήριο $0{,}000\,002$ μ. και ένας ιός $0{,}000\,000\,1$ μ. Ο Στέλιος προσπαθεί να τα συγκρίνει και χάνεται στα μηδενικά: «Ποιο είναι μικρότερο, το $0{,}000\,007$ ή το $0{,}000\,002$;».

Η Πολυξένη θυμάται ότι για τους πολύ μεγάλους αριθμούς η κυρία Αλεξάνδρα τους είχε δείξει να τους γράφουν με δυνάμεις του $10$: $1{,}496\cdot10^{8}$ χιλιόμετρα η απόσταση Γης -- Ήλιου. «Μπορούμε να κάνουμε το ίδιο και για τους πολύ μικρούς;» ρωτάει. «Τώρα που ξέρουμε τους αρνητικούς εκθέτες, ναι», απαντάει η καθηγήτρια. Το $0{,}000\,07$ γράφεται $7\cdot10^{-5}$, και ο ιός $1\cdot10^{-7}$.

Η Μυρτώ προσθέτει μια ιστορική σημείωση: «Αυτά τα μεγέθη μπορούν να τα μετρήσουν χάρη στα επιστημονικά όργανα και στη γλώσσα των αριθμών. Η Υπατία, η μεγάλη μαθηματικός και αστρονόμος της Αλεξάνδρειας (περίπου 370--415 μ.Χ.), είχε δουλέψει σε όργανα μέτρησης του ουρανού, τον αστρολάβο». Η τυποποιημένη μορφή φέρνει τους πολύ μεγάλους και τους πολύ μικρούς αριθμούς στο ίδιο τετράδιο, με την ίδια γλώσσα.
\end{anagke}

\begin{lexi}{τυποποιημένη μορφή, τάξη μεγέθους}{τυποποιημένη μορφή, τάξη μεγέθους}
Όπως είδαμε, \lexiword{τυποποιημένη} είναι η μορφή που έχει το ίδιο «τύπο» για όλους. Η \lexiword{τάξη μεγέθους} ενός αριθμού είναι ο εκθέτης του $10$ στην τυποποιημένη του μορφή: μας λέει «πόσο μεγάλος ή πόσο μικρός» είναι, με μια ματιά.
\end{lexi}

\begin{metafrasi}
Το πάχος της τρίχας. & $0{,}000\,07=7\cdot10^{-5}$ μ.\\
Το αιμοσφαίριο. & $0{,}000\,007=7\cdot10^{-6}$ μ.\\
Ο ιός. & $0{,}000\,000\,1=1\cdot10^{-7}$ μ.\\
Η ταχύτητα του φωτός. & $300\,000\,000=3\cdot10^{8}$ μ./s\\
Η οθόνη του κομπιουτεράκι. & $4.5\mathrm{E}{-}4=4{,}5\cdot10^{-4}$\\
\end{metafrasi}

\tmima{Η θεωρία}

\begin{orismos}{τυποποιημένη μορφή}
Κάθε θετικός ρητός μπορεί να γραφεί στη μορφή $\alpha\cdot10^{\nu}$ με $1\le\alpha<10$ και $\nu$ \textbf{ακέραιο}. Αυτή λέγεται \textbf{τυποποιημένη μορφή}. Για τους μεγάλους αριθμούς ο $\nu$ είναι θετικός, για τους μικρούς ($<1$) αρνητικός.
\end{orismos}

\begin{kanonas}{τυποποίηση}
\begin{itemize}[leftmargin=1.2em,itemsep=0.3ex,topsep=2pt]
\item Μεταφέρω την υποδιαστολή ώστε να μείνει ένα μη μηδενικό ψηφίο πριν από αυτήν και βρίσκω τον $\alpha$.
\item Ο εκθέτης $\nu$ είναι το πλήθος των θέσεων που μετακίνησα, με πρόσημο $+$ αν μετακίνησα την υποδιαστολή αριστερά (αριθμός μεγάλος) και $-$ αν τη μετακίνησα δεξιά (αριθμός μικρός).
\end{itemize}
\end{kanonas}

Έτσι $0{,}00045$: μετακινώ $4$ θέσεις δεξιά και παίρνω $4{,}5$: $0{,}00045=4{,}5\cdot10^{-4}$. Και $3\,420\,000=3{,}42\cdot10^{6}$. Για να γυρίσω πίσω, πολλαπλασιάζω τον $\alpha$ με το $10^\nu$: $3{,}2\cdot10^{-5}=0{,}000032$.

\tmimasmall{Σύγκριση και πράξεις}

Συγκρίνω πρώτα τους εκθέτες: μεγαλύτερος εκθέτης σημαίνει μεγαλύτερο αριθμό ($4\cdot10^{-5}>7\cdot10^{-6}$ αφού $-5>-6$). Αν οι εκθέτες είναι ίσοι συγκρίνω τους $\alpha$. Για πράξεις: στον πολλαπλασιασμό πολλαπλασιάζω τους $\alpha$ και προσθέτω τους εκθέτες, στη διαίρεση διαιρώ τους $\alpha$ και αφαιρώ τους εκθέτες, και μετά ξαναγράφω σε τυποποιημένη μορφή αν χρειάζεται:
$$(3\cdot10^{-3})(2\cdot10^{5})=6\cdot10^{2},\qquad(6\cdot10^{8}):(2\cdot10^{-2})=3\cdot10^{10},\qquad(5\cdot10^{4})(4\cdot10^{-2})=20\cdot10^{2}=2\cdot10^{3}.$$

\begin{prosoxi}
\begin{itemize}[leftmargin=1.2em,itemsep=0.3ex,topsep=2pt]
\item Ο $\alpha$ είναι πάντα από $1$ έως λιγότερο από $10$. Το $12\cdot10^{-4}$ δεν είναι τυποποιημένο: είναι $1{,}2\cdot10^{-3}$.
\item Στους μικρούς αριθμούς ο εκθέτης είναι \emph{αρνητικός}: το $0{,}00045$ δεν είναι $4{,}5\cdot10^{4}$ αλλά $4{,}5\cdot10^{-4}$.
\item Ο εκθέτης $-4$ δεν σημαίνει «τέσσερα μηδενικά» αλλά «η υποδιαστολή πήγε τέσσερις θέσεις δεξιά»: μετά την υποδιαστολή υπάρχουν $3$ μηδενικά και μετά το $4$.
\item Στη σύγκριση: $10^{-5}$ είναι μεγαλύτερο από $10^{-6}$. Όσο πιο αρνητικός ο εκθέτης, τόσο μικρότερος ο αριθμός.
\end{itemize}
\end{prosoxi}

\tmima{Λυμένα παραδείγματα}

\paradeigma{Μικροί αριθμοί}
\tixero Οι $0{,}00045$, $0{,}000\,000\,7$, $0{,}0302$.\par
\tithelo Να τους γράψω σε τυποποιημένη μορφή.\par
\possindeo Μετακινώ την υποδιαστολή δεξιά μέχρι ένα μη μηδενικό ψηφίο πριν από αυτήν. $0{,}00045$: $4$ θέσεις, $4{,}5\cdot10^{-4}$. $0{,}000\,000\,7$: $7$ θέσεις, $7\cdot10^{-7}$. $0{,}0302$: $2$ θέσεις, $3{,}02\cdot10^{-2}$.\par
\apantisi $4{,}5\cdot10^{-4}$, $7\cdot10^{-7}$, $3{,}02\cdot10^{-2}$.\par
\elegxos $4{,}5\cdot10^{-4}=\frac{4{,}5}{10\,000}=0{,}00045$.

\paradeigma{Από τυποποιημένη σε δεκαδική}
\tixero Οι $3{,}2\cdot10^{-5}$, $9{,}1\cdot10^{-3}$, $2{,}5\cdot10^{6}$.\par
\tithelo Να τους γράψω κανονικά.\par
\possindeo Αρνητικός εκθέτης: υποδιαστολή αριστερά. $3{,}2\cdot10^{-5}=0{,}000032$. $9{,}1\cdot10^{-3}=0{,}0091$. Θετικός εκθέτης: υποδιαστολή δεξιά. $2{,}5\cdot10^6=2\,500\,000$.\par
\apantisi $0{,}000032$, $0{,}0091$, $2\,500\,000$.\par
\elegxos $3{,}2\cdot10^{-5}$: $5$ θέσεις αριστερά από το $3$: $0{,}000032$.

\paradeigma{Σύγκριση μεγεθών}
\tixero Η τρίχα $7\cdot10^{-5}$ μ., το αιμοσφαίριο $7\cdot10^{-6}$ μ., το βακτήριο $2\cdot10^{-6}$ μ., ο ιός $1\cdot10^{-7}$ μ.\par
\tithelo Να τα βάλω σε σειρά από το μικρότερο και να βρω πόσες φορές η τρίχα είναι χοντρότερη από το αιμοσφαίριο.\par
\possindeo Συγκρίνω εκθέτες: $-7<-6<-5$. Ο ιός ο μικρότερος. Μετά τα δύο με εκθέτη $-6$: $2<7$, δηλαδή το βακτήριο και μετά το αιμοσφαίριο. Τελευταία η τρίχα. Ο λόγος: $\frac{7\cdot10^{-5}}{7\cdot10^{-6}}=10^{-5-(-6)}=10$.\par
\apantisi $1\cdot10^{-7}<2\cdot10^{-6}<7\cdot10^{-6}<7\cdot10^{-5}$. Η τρίχα είναι $10$ φορές χοντρότερη.\par
\elegxos $0{,}00007=10\cdot0{,}000007$.

\paradeigma{Πράξεις}
\tixero Οι $(3\cdot10^{-3})\cdot(2\cdot10^{5})$, $(6\cdot10^{8}):(2\cdot10^{-2})$, $(5\cdot10^{4})\cdot(4\cdot10^{-2})$.\par
\tithelo Να τις υπολογίσω σε τυποποιημένη μορφή.\par
\possindeo $3\cdot2=6$ και $-3+5=2$: $6\cdot10^2$. $6:2=3$ και $8-(-2)=10$: $3\cdot10^{10}$. $5\cdot4=20$ και $4-2=2$: $20\cdot10^2=2\cdot10^3$.\par
\apantisi $6\cdot10^{2}$, $3\cdot10^{10}$, $2\cdot10^{3}$.\par
\elegxos Το τρίτο: $50\,000\cdot0{,}04=2\,000$.

\paradeigma{Το φως του Ήλιου}
\tixero Το φως ταξιδεύει με $3\cdot10^{8}$ μ./s και ο Ήλιος απέχει $1{,}5\cdot10^{11}$ μ.\par
\tithelo Σε πόσο χρόνο φτάνει το φως στη Γη.\par
\possindeo $\frac{1{,}5\cdot10^{11}}{3\cdot10^{8}}=0{,}5\cdot10^{3}=5\cdot10^{2}$ s.\par
\apantisi $500$ s, περίπου $8{,}3$ λεπτά.\par
\elegxos $500\cdot3\cdot10^8=1{,}5\cdot10^{11}$.

\tmima{Ασκήσεις}

\zoni{Το κατάλαβα}{μικροί και μεγάλοι αριθμοί}

\ex Γράψε σε τυποποιημένη μορφή: α) $0{,}0005$ \quad β) $0{,}000\,072$ \quad γ) $0{,}3$ \quad δ) $45\,000$ \quad ε) $0{,}000\,000\,8$ \quad στ) $6\,300\,000$.

\ex Γράψε σε δεκαδική μορφή: α) $3{,}2\cdot10^{-5}$ \quad β) $9{,}1\cdot10^{-3}$ \quad γ) $4\cdot10^{-1}$ \quad δ) $2{,}5\cdot10^{6}$ \quad ε) $1{,}01\cdot10^{-4}$ \quad στ) $7\cdot10^{0}$.

\ex Γράψε σωστά σε τυποποιημένη μορφή: α) $0{,}3\cdot10^{-2}$ \quad β) $12\cdot10^{-4}$ \quad γ) $5{,}5\cdot10^{-3}$ \quad δ) $10\cdot10^{2}$.

\ex Σύγκρινε: α) $4\cdot10^{-5}$ και $7\cdot10^{-6}$ \quad β) $2{,}5\cdot10^{-3}$ και $2{,}5\cdot10^{-2}$ \quad γ) $6\cdot10^{4}$ και $9\cdot10^{-4}$ \quad δ) $3{,}1\cdot10^{-8}$ και $3{,}09\cdot10^{-8}$.

\ex Γράψε σε αύξουσα σειρά: $3\cdot10^{-3}$, $5\cdot10^{-5}$, $2\cdot10^{2}$, $9\cdot10^{-1}$, $1{,}2\cdot10^{-3}$.

\ex Υπολόγισε σε τυποποιημένη μορφή: α) $(3\cdot10^{-3})(2\cdot10^{5})$ \quad β) $(6\cdot10^{8}):(2\cdot10^{-2})$ \quad γ) $(5\cdot10^{4})(4\cdot10^{-2})$ \quad δ) $(8\cdot10^{-6}):(4\cdot10^{-3})$.

\ex Ο υπολογιστής δείχνει $4.5\mathrm{E}{-}4$, $2.5\mathrm{E}7$ και $7\mathrm{E}{-}9$. Ποιοι αριθμοί είναι;

\ex Σωστό ή λάθος; α) $0{,}004=4\cdot10^{-3}$. \quad β) Το $12\cdot10^{-4}$ είναι τυποποιημένο. \quad γ) Το $10^{-5}$ είναι μικρότερο από το $10^{-6}$. \quad δ) $2\cdot10^{-3}=0{,}002$. \quad ε) Όσο μεγαλύτερος ο εκθέτης, τόσο μεγαλύτερος ο αριθμός. \quad στ) $5{,}5\cdot10^{-3}>5{,}5\cdot10^{-2}$.

\zoni{Το χρησιμοποιώ}{ο μικρόκοσμος και ο μακρόκοσμος}

\ex \textbf{Ο μικρόκοσμος.} Τρίχα $7\cdot10^{-5}$ μ., αιμοσφαίριο $7\cdot10^{-6}$ μ., βακτήριο $2\cdot10^{-6}$ μ., ιός $10^{-7}$ μ. Βάλε τα σε αύξουσα σειρά και βρες πόσες φορές η τρίχα είναι χοντρότερη από το αιμοσφαίριο.

\ex \textbf{Το ηλεκτρόνιο.} Η μάζα του ηλεκτρονίου είναι $9{,}1\cdot10^{-31}$ kg. Πόσα μηδενικά υπάρχουν μετά την υποδιαστολή πριν από το $9$ αν τη γράψεις δεκαδικά;

\ex \textbf{Το φως.} Το φως ταξιδεύει με $3\cdot10^{8}$ μ./s και ο Ήλιος απέχει από τη Γη $1{,}5\cdot10^{11}$ μ. Σε πόσο χρόνο φτάνει το φως στη Γη;

\ex \textbf{Το DNA.} Το πλάτος του DNA είναι $2$ nm και $1\ \mathrm{nm}=10^{-9}$ μ. Γράψε το σε μέτρα σε τυποποιημένη μορφή.

\ex \textbf{Το χαρτί.} Ένα φύλλο χαρτιού έχει πάχος $10^{-4}$ μ. Πόσο είναι το πάχος $500$ φύλλων σε μέτρα και σε εκατοστά;

\ex \textbf{Η άμμος.} Ένας κόκκος άμμου έχει μάζα $10^{-6}$ kg. Πόσοι κόκκοι ζυγίζουν $1$ kg;

\zoni{Το σκέφτομαι}{ο εκθέτης που μικραίνει}

\ex Γιατί ο εκθέτης είναι αρνητικός για αριθμούς μικρότερους από το $1$;

\ex Ποια είναι η τυποποιημένη μορφή του $1$; Και του $10$;

\ex Σύγκρινε τα $2\cdot10^{-3}$ και $5\cdot10^{-4}$ με δύο τρόπους (δεκαδικά και εκθέτες).

\ex Αν πολλαπλασιάσεις το $3\cdot10^{-4}$ με το $10^{6}$ τι παίρνεις;

\ex Γιατί το $12\cdot10^{-4}$ δεν είναι τυποποιημένο; Γράψε το σωστά.

\zoni{Για το σπίτι}{ψάξε τα μεγέθη}

\ex Βρες το μέγεθος ενός ατόμου, ενός ιού και ενός βακτηρίου από αξιόπιστη πηγή και γράψε τα σε τυποποιημένη μορφή και σε αύξουσα σειρά.

\epanalipsi{7}

\tmima{Λεξικό του κεφαλαίου}

\renewcommand{\arraystretch}{1.5}
\begin{tabular*}{\linewidth}{@{\extracolsep{\fill}} >{\raggedright\arraybackslash}p{0.28\linewidth} >{\raggedright\arraybackslash}p{0.62\linewidth}}
\toprule
\textbf{Λέξη ή τύπος} & \textbf{Με μία πρόταση}\\
\midrule
ρητοί & Φυσικοί, κλάσματα και δεκαδικοί μαζί με τους αντίστοιχους αρνητικούς.\\
$|\alpha|$ & Η απόσταση του $\alpha$ από το $0$ στον άξονα: ποτέ αρνητική.\\
αντίθετοι & Ίδια απόλυτη τιμή, ετερόσημοι: $\alpha$ και $-\alpha$, με άθροισμα $0$.\\
σύγκριση & Μεγαλύτερος είναι ο δεξιότερος στον άξονα.\\
$\alpha-\beta$ & $=\alpha+(-\beta)$: προσθέτω τον αντίθετο.\\
γινόμενο, πηλίκο & Ομόσημοι: θετικό. Ετερόσημοι: αρνητικό. Όχι διαίρεση με το $0$.\\
περιοδικός & Δεκαδικός με ψηφία που επαναλαμβάνονται, π.χ. $\frac{1}{3}=0{,}\overline{3}$.\\
$\alpha^{\nu}$, $\alpha^0$, $\alpha^{-\nu}$ & $\alpha\cdots\alpha$, $1$ και $\frac{1}{\alpha^\nu}$.\\
$\alpha\cdot10^{\nu}$ & Τυποποιημένη μορφή: $1\le\alpha<10$, $\nu$ ακέραιος.\\
\bottomrule
\end{tabular*}

\tmima{Ασκήσεις επανάληψης}

\ex Εξέφρασε με ρητό αριθμό: α) $8$ βαθμοί κάτω από το μηδέν \quad β) χρέος $25$ € \quad γ) βάθος $40$ μ. κάτω από την επιφάνεια της θάλασσας.

\ex Ποιοι από τους $-3$, $\frac{4}{2}$, $0{,}5$, $-\frac{7}{7}$, $0$ είναι ακέραιοι;

\ex Βρες: α) $|-7{,}5|$ \quad β) $\left|\frac{3}{4}\right|$ \quad γ) τον αντίθετο του $-\frac{2}{3}$ \quad δ) τους $\alpha$ με $|\alpha|=4$.

\ex Γράψε σε αύξουσα σειρά: $-2{,}5$, $0$, $-\frac{1}{2}$, $3$, $-3$.

\ex Υπολόγισε: α) $(-7)+(+12)$ \quad β) $(-3{,}5)+(-1{,}5)$ \quad γ) $\left(-\frac{1}{2}\right)+\left(+\frac{1}{3}\right)$ \quad δ) $(+8)+(-8)$.

\ex Υπολόγισε: α) $5-(-3)$ \quad β) $(-4)-(+6)$ \quad γ) $(-2{,}5)-(-2{,}5)$ \quad δ) $\left(-\frac{1}{2}\right)-\frac{1}{4}$.

\ex Υπολόγισε το αλγεβρικό άθροισμα $-12+7-(-5)-(3-9)$.

\ex Υπολόγισε: α) $(-6)\cdot(-7)$ \quad β) $(+4)\cdot(-2{,}5)$ \quad γ) $\left(-\frac{1}{2}\right)(-4)(-3)$ \quad δ) $(-2)\cdot(5-8)$.

\ex Υπολόγισε: α) $(-48):(-6)$ \quad β) $15:(-5)$ \quad γ) $\left(-\frac{3}{4}\right):\left(-\frac{9}{2}\right)$ \quad δ) $0:(-7)$.

\ex Γράψε σε δεκαδική μορφή τα $\frac{5}{6}$ και $\frac{7}{20}$ και ως κλάσμα τον $0{,}\overline{18}$.

\ex Υπολόγισε: α) $(-2)^5$ \quad β) $-2^4$ \quad γ) $\left(\frac{1}{3}\right)^{-2}$ \quad δ) $5^0+2^{-3}$.

\ex Υπολόγισε με ιδιότητες: α) $2^5\cdot2^{-3}$ \quad β) $(3^{-1})^2$ \quad γ) $10^4:10^6$.

\ex Γράψε σε τυποποιημένη μορφή τα $0{,}00032$ και $45\,000\,000$, και υπολόγισε $(4\cdot10^{3})(5\cdot10^{-5})$.

\ex \textbf{Η πτώση.} Η θερμοκρασία πέφτει από $+4$° σε $-5$° σε $3$ ώρες. Ποιος είναι ο ρυθμός μεταβολής ανά ώρα;

\ex \textbf{Ερωτήσεις αυτοαξιολόγησης (πράξεις).} Σωστό ή λάθος; α) $|-5|=-5$. \quad β) $(-3)+(-4)=-7$. \quad γ) $5-(-2)=3$. \quad δ) $(-2)(-3)(-4)=-24$. \quad ε) $12:(-4)=3$. \quad στ) Ο αντίθετος του $0$ είναι το $0$.

\ex \textbf{Ερωτήσεις αυτοαξιολόγησης (δυνάμεις, δεκαδικοί).} Σωστό ή λάθος; α) $(-2)^0=1$. \quad β) $2^{-1}=-2$. \quad γ) Το $\frac{1}{7}$ είναι περιοδικός δεκαδικός. \quad δ) $10^{-3}>10^{-4}$. \quad ε) $(\alpha+\beta)^2=\alpha^2+\beta^2$. \quad στ) $3{,}2\cdot10^{-2}=0{,}032$.

\tmima{Ο χάρτης της επανάληψης}

Μετά τις λύσεις, δες σε ποιες ασκήσεις δυσκολεύτηκες. Ο πίνακας σου λέει ποια παράγραφο να ξαναδιαβάσεις.

\renewcommand{\arraystretch}{1.5}
\begin{tabular*}{\linewidth}{@{\extracolsep{\fill}} >{\raggedright\arraybackslash}p{0.38\linewidth} >{\raggedright\arraybackslash}p{0.52\linewidth}}
\toprule
\textbf{Αν δυσκολεύτηκες στις} & \textbf{ξαναδιάβασε τις παραγράφους}\\
\midrule
Ε7.1, Ε7.2 & 7.1 (ρητοί, άξονας)\\
Ε7.3, Ε7.4 & 7.2 (απόλυτη τιμή, αντίθετοι, σύγκριση)\\
Ε7.5 & 7.3 (πρόσθεση)\\
Ε7.6, Ε7.7 & 7.4 (αφαίρεση, παρενθέσεις)\\
Ε7.8 & 7.5 (πολλαπλασιασμός)\\
Ε7.9, Ε7.14 & 7.6 (διαίρεση)\\
Ε7.10 & 7.7 (δεκαδική μορφή)\\
Ε7.11, Ε7.12 & 7.8 και 7.9 (δυνάμεις)\\
Ε7.13 & 7.10 (τυποποιημένη μορφή)\\
Ε7.15, Ε7.16 & 7.1 έως 7.10\\
\bottomrule
\end{tabular*}

